\documentclass[12pt, a4paper]{article}

% ── Packages ──────────────────────────────────────────────────────────────────
\usepackage[a4paper, margin=2.5cm]{geometry}
\usepackage{graphicx}
\usepackage{float}
\usepackage{amsmath}
\usepackage{amssymb}
\usepackage{xcolor}
\usepackage{mdframed}
\usepackage{parskip}
\usepackage{enumitem}
\usepackage{titlesec}
\usepackage{booktabs}
\usepackage{hyperref}
\usepackage[T1]{fontenc}

% ── Graphics path ─────────────────────────────────────────────────────────────
\graphicspath{{../../Images_for_latex/Lecture_7/}}

% ── Colours ───────────────────────────────────────────────────────────────────
\definecolor{keyblue}{RGB}{30, 80, 160}
\definecolor{keybluebg}{RGB}{220, 235, 255}
\definecolor{noteorange}{RGB}{200, 100, 0}
\definecolor{noteorangebg}{RGB}{255, 240, 215}
\definecolor{warnred}{RGB}{180, 30, 30}
\definecolor{warnredbg}{RGB}{255, 220, 220}
\definecolor{defgreen}{RGB}{20, 110, 50}
\definecolor{defgreenbg}{RGB}{215, 245, 225}
\definecolor{headernavy}{RGB}{25, 35, 90}

% ── Box environments ──────────────────────────────────────────────────────────
\newmdenv[linecolor=keyblue,backgroundcolor=keybluebg,linewidth=1.5pt,
  leftline=true,rightline=false,topline=false,bottomline=false,
  innerleftmargin=10pt,innerrightmargin=8pt,
  innertopmargin=6pt,innerbottommargin=6pt,
  skipabove=8pt,skipbelow=8pt]{keybox}

\newmdenv[linecolor=noteorange,backgroundcolor=noteorangebg,linewidth=1.5pt,
  leftline=true,rightline=false,topline=false,bottomline=false,
  innerleftmargin=10pt,innerrightmargin=8pt,
  innertopmargin=6pt,innerbottommargin=6pt,
  skipabove=8pt,skipbelow=8pt]{notebox}

\newmdenv[linecolor=warnred,backgroundcolor=warnredbg,linewidth=1.5pt,
  leftline=true,rightline=false,topline=false,bottomline=false,
  innerleftmargin=10pt,innerrightmargin=8pt,
  innertopmargin=6pt,innerbottommargin=6pt,
  skipabove=8pt,skipbelow=8pt]{warnbox}

\newmdenv[linecolor=defgreen,backgroundcolor=defgreenbg,linewidth=1.5pt,
  leftline=true,rightline=false,topline=false,bottomline=false,
  innerleftmargin=10pt,innerrightmargin=8pt,
  innertopmargin=6pt,innerbottommargin=6pt,
  skipabove=8pt,skipbelow=8pt]{defbox}

% ── Section formatting ────────────────────────────────────────────────────────
\titleformat{\section}[block]{\large\bfseries\color{headernavy}}{}{0pt}{}
  [\vspace{2pt}\rule{\linewidth}{0.8pt}\vspace{4pt}]
\titleformat{\subsection}[block]{\normalsize\bfseries\color{keyblue}}{}{0pt}{}

% ── Document ──────────────────────────────────────────────────────────────────
\begin{document}

\begin{center}
  {\LARGE\bfseries\color{headernavy} Sensors \& Systems}\\[4pt]
  {\large Ranging-based Sensors}\\[2pt]
  {\normalsize Jan Bendtsen, Torben Knudsen \quad|\quad Electronic Systems, 8th Semester}
\end{center}

\vspace{4pt}\hrule height 1.2pt\vspace{12pt}

% ==============================================================================
\section*{Section 1 \quad Ranging-based Sensors \& Applications}
% ==============================================================================

A \textbf{ranging-based sensor} measures the \emph{distance} from the sensor to
a target or surface.
Rather than capturing an image or measuring a physical quantity at the sensor's
own location, it sends out some kind of wave or pulse, waits for the echo to
return, and infers distance from the round-trip travel time or phase.
The three dominant technologies are \textbf{Sonar}, \textbf{Radar}, and
\textbf{LiDAR} --- each using a different part of the wave spectrum.

\begin{figure}[H]
  \centering
  \includegraphics[width=0.82\textwidth]{slide-03.jpg}
  \caption{Overview of the three main ranging sensor families: Radar (radio
    waves, long range, penetrates weather), LiDAR (laser light, high spatial
    resolution), and Sonar (sound waves, primarily underwater).}
\end{figure}

% ------------------------------------------------------------------------------
\subsection*{1.1 \quad Sonar --- Sound Navigation and Ranging}
% ------------------------------------------------------------------------------

\textbf{Sonar} uses \emph{sound propagation} to navigate and measure distances.
It is most commonly used \emph{underwater} (submarines, ship navigation, fish
finding), where radio waves cannot propagate effectively.

\begin{figure}[H]
  \centering
  \includegraphics[width=0.72\textwidth]{slide-04.jpg}
  \caption{Active sonar (left) emits a sound pulse and listens for echoes
    returning from objects.  Passive sonar (right) only listens for sounds
    emitted by other vessels or targets.}
\end{figure}

There are two operating modes:

\begin{itemize}[leftmargin=1.4em]
  \item \textbf{Passive sonar} --- simply listens for sounds emitted by other
    objects (e.g.\ an enemy submarine's engine noise).  No pulse is transmitted,
    so the sensor itself remains undetected.
  \item \textbf{Active sonar} --- emits a pulse of sound and listens for the
    returning echo.  Allows precise distance measurement but reveals the
    presence of the sensor.
\end{itemize}

\begin{notebox}
  \textbf{Nature uses sonar too.}
  Bats, toothed whales (e.g.\ dolphins), and certain shrews use \emph{active}
  biological sonar (echo-location) to navigate and hunt.
  Dolphins can emit up to 600 clicks per second.
  The focused beam is shaped by a fatty organ called the \emph{melon}.
\end{notebox}

% ------------------------------------------------------------------------------
\subsection*{1.2 \quad Radar --- Radio Detection and Ranging}
% ------------------------------------------------------------------------------

\textbf{Radar} uses \emph{radio waves} to determine the distance, direction
(azimuth and elevation angles), and \emph{radial velocity} of objects relative
to the sensor site.

\begin{figure}[H]
  \centering
  \includegraphics[width=0.62\textwidth]{slide-05.jpg}
  \caption{A traditional rotating-dish radar antenna.  Modern high-tech radar
    systems use digital signal processing and machine learning to extract
    useful information even at very high noise levels.}
\end{figure}

Key properties of Radar:
\begin{itemize}[leftmargin=1.4em]
  \item Uses radio waves, typically in the range \textbf{30\,cm -- 3\,mm}
    wavelength (microwave band).
  \item Penetrates \textbf{fog, clouds, rain, and darkness} --- making it ideal
    for aviation, maritime, and automotive applications where LiDAR would
    struggle in bad weather.
  \item Can also measure \textbf{radial velocity} directly via the Doppler
    effect (frequency shift of the returning wave).
  \item Used to detect and track aircraft, ships, spacecraft, guided missiles,
    motor vehicles, weather formations, and terrain.
\end{itemize}

\begin{notebox}
  \textbf{What do azimuth, elevation, and radial velocity actually mean?}

  A radar measures three geometric quantities about a target simultaneously:

  \begin{itemize}[leftmargin=1.2em,itemsep=3pt]
    \item \textbf{Azimuth} --- the \emph{horizontal} angle to the target,
      like a compass bearing.  Standing at the radar: 0$^\circ$ = North, 90$^\circ$ = East,
      180$^\circ$ = South, and so on.  The rotating dish sweeps through all azimuth
      angles to scan the full horizon.

    \item \textbf{Elevation} --- the \emph{vertical} angle above the horizon.
      0$^\circ$ is flat along the ground, 90$^\circ$ is straight up.  Together, azimuth +
      elevation give the full 3D \emph{direction} to the target (like
      spherical coordinates).

    \item \textbf{Radial velocity} --- the component of the target's velocity
      that is directed \emph{directly toward or away from the radar}.  This is
      measured via the \textbf{Doppler effect}: a target moving toward the
      radar compresses the returning radio wave (higher frequency); a target
      moving away stretches it (lower frequency).  The frequency shift is
      directly proportional to the radial velocity.

      \textit{Important:} this is \emph{not} the full velocity vector.  A car
      driving sideways past the radar has nearly zero radial velocity even if
      it is travelling at 100\,km/h --- the beam sees it moving neither toward
      nor away.
  \end{itemize}

  Together: \textbf{azimuth + elevation + range} give the full 3D
  \emph{position} of the target; \textbf{radial velocity} gives how fast it
  is approaching or receding.  This is why radar is so effective for tracking
  aircraft: one measurement gives both position and approach speed.
\end{notebox}

\begin{notebox}
  \textbf{Active vs.\ passive radar --- and how range is actually computed}

  Just like sonar, radar can operate in two modes:
  \begin{itemize}[leftmargin=1.2em,itemsep=3pt]
    \item \textbf{Active radar} --- transmits its own radio pulse and listens
      for the returning echo.  The classic rotating-dish radar is active.
    \item \textbf{Passive radar} --- transmits \emph{nothing}.  Instead it
      exploits existing radio signals in the environment --- FM radio
      stations, TV broadcasts, WiFi, mobile towers --- as
      \textit{illuminators of opportunity}, and detects their reflections
      off targets.  Because no pulse is emitted, the radar itself stays
      completely undetectable; similar in spirit to passive sonar.
  \end{itemize}

  \textbf{How range is computed} depends on the radar type:
  \begin{itemize}[leftmargin=1.2em,itemsep=3pt]
    \item \textbf{Pulsed radar} --- emits a short burst and measures the
      round-trip \emph{time of flight}.  Same principle as pulsed LiDAR:
      \[
        R = \tfrac{1}{2}\,c\,t
      \]
      where $c$ is the speed of light and $t$ is the measured delay.

    \item \textbf{FMCW radar} (Frequency-Modulated Continuous Wave) ---
      transmits a continuous signal whose frequency is swept linearly up and
      down in a ramp (a ``chirp'').  Because the received echo is delayed
      relative to the transmitted signal, there is a constant
      \emph{frequency difference} $\Delta f$ between what is being
      transmitted and what is being received at any instant.  This
      frequency difference is directly proportional to range:
      \[
        R \propto \Delta f
      \]
      FMCW is widely used in automotive radar because it can measure range
      and radial velocity \emph{simultaneously} from the same waveform.

    \item \textbf{Doppler shift} --- the \emph{change} in frequency of the
      returned wave (not the difference described above) gives the radial
      velocity, not the range directly.
  \end{itemize}

  In summary: \textbf{time of flight} $\to$ range (pulsed);
  \textbf{frequency difference} $\to$ range (FMCW);
  \textbf{frequency shift} $\to$ radial velocity (Doppler).
\end{notebox}

% ------------------------------------------------------------------------------
\subsection*{1.3 \quad LiDAR --- Light Detection and Ranging}
% ------------------------------------------------------------------------------

\textbf{LiDAR} uses \emph{laser light} (ultraviolet, visible, or near-infrared)
to measure distances with very high spatial resolution.

\begin{figure}[H]
  \centering
  \includegraphics[width=0.72\textwidth]{slide-06.jpg}
  \caption{LiDAR point cloud of two cars.  A spinning LiDAR head sweeps laser
    beams in all directions; each reflection gives one 3D data point, building
    up a dense point cloud of the environment.}
\end{figure}

Key properties of LiDAR:
\begin{itemize}[leftmargin=1.4em]
  \item Most systems work in \textbf{pulsed} ranging mode: a very short pulse
    ($<50$\,ns) is emitted, and part of the reflected pulse returns to the
    sensor.  Distance is computed from the time of flight.
  \item Unlike radar, LiDAR uses \textbf{laser light} (typically 903 or 905\,nm
    near-infrared), giving \emph{much finer angular resolution} but poorer
    performance in rain, fog, or dust.
  \item Due to the narrow beam, LiDAR uses \textbf{swivel mechanisms or phased
    arrays} to scan across a field of view, producing a full 3D point cloud.
\end{itemize}

\begin{notebox}
  \textbf{What a returning laser pulse actually tells you --- range and colour}

  Every time the laser fires, two independent pieces of information come back
  from the returning pulse:

  \begin{itemize}[leftmargin=1.2em,itemsep=3pt]
    \item \textbf{When} it returns $\;\to\;$ \textbf{range.}
      The round-trip time of flight $t$ gives the distance directly:
      \[
        R = \tfrac{1}{2}\,c\,t
      \]
      This is exactly the same principle as pulsed radar and pulsed sonar ---
      only the wave speed $c$ differs.

    \item \textbf{How strongly} it returns $\;\to\;$ \textbf{reflectivity
      (intensity / ``colour'').}
      The \emph{amplitude} of the returned pulse tells you how much laser
      energy the surface reflected back toward the sensor.
      A highly reflective surface (white road markings, metal signs,
      retroreflectors) returns a strong pulse $\to$ bright/high-intensity
      point.
      A dark or absorbing surface (black asphalt, rubber tyres) returns a weak
      pulse $\to$ dim/low-intensity point.
      This is why point clouds can be \emph{coloured by return intensity}
      --- each point carries both its 3D position and a reflectivity value.
      The coloured-intensity image from slide~19 is exactly this: brighter
      pixels are more reflective surfaces.
  \end{itemize}

  \textbf{Why wavelength (frequency) matters for material identification.}
  Different materials reflect different wavelengths differently.  LiDAR
  typically uses near-infrared (903--905\,nm) because:
  \begin{itemize}[leftmargin=1.2em,itemsep=2pt]
    \item It is eye-safe at reasonable power levels.
    \item Water and chlorophyll \emph{absorb} near-infrared, so vegetation
      gives a weaker return than dry surfaces at the same distance --- the
      sensor can start to distinguish material type from intensity alone.
    \item Using \emph{two wavelengths simultaneously} (multi-spectral LiDAR)
      gives two intensity values per point.  Because each material has its
      own reflectivity ``fingerprint'' across wavelengths, comparing the two
      intensities allows the system to \emph{classify} what the surface is
      made of (e.g.\ tree canopy vs.\ ground vs.\ building facade) even if
      they are at the same height.
  \end{itemize}

  In short: \textbf{time} $\to$ range;
  \textbf{amplitude} $\to$ reflectivity / point colour;
  \textbf{wavelength choice} $\to$ material sensitivity.
\end{notebox}

% ── Comparison table ──────────────────────────────────────────────────────────
\begin{notebox}
  \textbf{Sensor comparison at a glance}

  \vspace{4pt}
  \begin{tabular}{@{}lllll@{}}
    \toprule
    \textbf{Sensor} & \textbf{Wave type} & \textbf{Range} &
      \textbf{Resolution} & \textbf{Weather} \\
    \midrule
    Sonar  & Sound ($\sim$kHz--MHz)  & Medium (water) & Moderate   & N/A (water) \\
    Radar  & Radio (mm--cm)          & Long (km)       & Low--Med   & Excellent \\
    LiDAR  & Laser (near-IR)         & Short--Med      & Very high  & Poor (fog/rain) \\
    \bottomrule
  \end{tabular}
\end{notebox}

% ------------------------------------------------------------------------------
\clearpage
\subsection*{1.4 \quad Applications}
% ------------------------------------------------------------------------------

Ranging sensors are found wherever a machine needs to \emph{perceive} its
physical environment in 3D.  In some applications a \emph{single} sensor type
is sufficient; in others, \emph{combining} several sensors is necessary because
each has blind spots the others cover.

\textbf{Autonomous car navigation} --- sensor combination (slide~7):

\begin{figure}[H]
  \centering
  \includegraphics[width=0.78\textwidth]{slide-07.jpg}
  \caption{Three complementary sensors on an autonomous vehicle, each covering
    what the others cannot.}
\end{figure}

This is the clearest example of \emph{why} sensors are combined:
\begin{itemize}[leftmargin=1.4em,itemsep=3pt]
  \item \textbf{Radar alone} can detect that \emph{something} is ahead and
    how fast it is approaching (radial velocity via Doppler), and it works in
    fog, rain, and darkness.  But its spatial resolution is too low to tell
    whether that something is a pedestrian, a parked car, or a road sign.
  \item \textbf{LiDAR alone} builds a precise 3D point cloud of the
    surroundings (sub-centimetre resolution) and can identify shapes of objects
    well.  But it fails in heavy rain or fog because water droplets scatter
    the laser, and it cannot read text or colour (speed-limit signs, traffic
    lights).
  \item \textbf{Camera alone} reads colour, text, and rich visual detail
    (signs, lane markings, traffic lights), but gives no direct range
    measurement and struggles in darkness or glare.
\end{itemize}

Used \emph{together}: radar handles adverse weather and provides velocity;
LiDAR maps the 3D geometry of the scene in good conditions; the camera reads
visual information no ranging sensor can capture.  Each sensor fills in the
gaps of the others --- this is called \textbf{sensor fusion}.
\clearpage
\textbf{Underwater echo-location} --- sonar alone (slide~8):

\begin{figure}[H]
  \centering
  \includegraphics[width=0.72\textwidth]{slide-08.jpg}
  \caption{A dolphin using biological active sonar.  Sound pulses emitted from
    the head reflect off prey; the returning echoes are used to estimate
    distance and direction.  No combination needed --- sonar is the only
    practical ranging method underwater, since light and radio waves are
    absorbed within metres.}
\end{figure}

Here a \emph{single} sensor type is used because the medium (water) rules out
all alternatives: radio waves are absorbed almost immediately underwater, and
visible light scatters strongly beyond a few metres.  Sound, however,
propagates very efficiently in water.  Active sonar alone is therefore
sufficient to sense range and direction to targets.
\clearpage
\textbf{Geo-surveying} --- LiDAR alone (slide~9):

\begin{figure}[H]
  \centering
  \includegraphics[width=0.78\textwidth]{slide-09.jpg}
  \caption{Airborne LiDAR geo-survey.  The aircraft sweeps a laser beam
    downward; each reflected pulse gives one 3D ground point.  The result is a
    high-resolution digital terrain model used in surveying, archaeology,
    geology, forestry, and atmospheric physics.}
\end{figure}

In open-air surveying from an aircraft, LiDAR alone is sufficient: the goal is
purely to measure the 3D geometry of the terrain below with the highest
possible resolution.  Radar would give longer range but far lower spatial
resolution, and a camera gives no direct range.  LiDAR's ability to penetrate
gaps in tree canopy and return \emph{multiple echoes} (canopy return, lower
canopy, ground return --- see the waveform model in Section~3) makes it
uniquely suited to this task.


% ------------------------------------------------------------------------------
\clearpage
\subsection*{1.5 \quad Essential Components of a LiDAR Sensor}
% ------------------------------------------------------------------------------

Before diving into the measurement physics, it helps to understand the
\emph{hardware} that makes LiDAR work.

\begin{figure}[H]
  \centering
  \includegraphics[width=0.72\textwidth]{slide-10.jpg}
  \caption{Cut-away of a rotating LiDAR head showing the five key components:
    (1)~\textbf{Laser source} --- emits the outgoing pulse.
    (2)~\textbf{Tilting mirror} --- steers the beam vertically.
    (3)~\textbf{Servo motor} --- rotates the entire head horizontally for full
    360$^\circ$ azimuth coverage.
    (4)~\textbf{Optical rotary encoders} --- record the exact azimuth and
    elevation angles at the moment of each pulse, so every returned distance
    can be assigned its correct angular direction.
    (5)~\textbf{Receiver} --- captures the returning laser pulse.}
\end{figure}

\begin{itemize}[leftmargin=1.4em]
  \item The \textbf{laser source} emits a short, powerful pulse of coherent
    light.  The pulse duration directly affects range resolution
    (shorter pulse $\Rightarrow$ finer resolution --- see Section~2).
  \item The \textbf{tilting mirror + servo motor} combination allows the sensor
    to build a full 3D scan: the motor rotates the head through 360$^\circ$ in azimuth
    while the mirror tilts through a vertical field of view.
  \item The \textbf{optical rotary encoders} are critical for accuracy ---
    without knowing the exact pointing angle at the moment of each pulse, the
    3D coordinates of the return point cannot be computed.
  \item The \textbf{receiver} (photodetector) converts the returning photons
    back into an electrical signal.  The time between transmit and receive is
    measured to give the distance.
\end{itemize}

% ── Key takeaways ─────────────────────────────────────────────────────────────
\begin{keybox}
  \textbf{Section 1 --- Key Takeaways}
  \begin{itemize}[leftmargin=1.2em,itemsep=2pt]
    \item Ranging sensors measure distance by sending a wave and timing (or
      phase-comparing) the return.  The three main families are
      \textbf{Sonar} (sound), \textbf{Radar} (radio waves), and
      \textbf{LiDAR} (laser light).
    \item \textbf{Sonar} is optimal underwater; active mode measures range,
      passive mode only listens.
    \item \textbf{Radar} has the longest range and works in all weather, but
      lower spatial resolution.  It can also measure radial velocity via the
      Doppler effect.
    \item \textbf{LiDAR} gives the highest spatial resolution (sub-centimetre
      point clouds) but requires clear air.  It is rapidly becoming the
      dominant sensor in engineering applications (robotics, autonomous
      vehicles, mapping).
    \item A rotating LiDAR head combines a laser source, steering mirror,
      servo motor, rotary encoders (for angle), and a receiver.  The encoders
      are what turn a 1D range into a 3D point.
  \end{itemize}
\end{keybox}

% ==============================================================================
\newpage
\section*{Section 2 \quad Measurement Principles: Pulsed \& Phased LiDAR}
% ==============================================================================

Now that we know \emph{what} a LiDAR is made of, we look at \emph{how} it
actually measures distance.
There are two fundamentally different approaches: \textbf{pulsed} LiDAR, which
uses time of flight, and \textbf{phased} (continuous-wave) LiDAR, which uses
the phase of a modulated wave.
Both ultimately rely on the finite speed of light, but they extract the
distance in different ways and come with different trade-offs in range and
precision.

% ------------------------------------------------------------------------------
\subsection*{2.1 \quad Pulsed LiDAR --- Time of Flight}
% ------------------------------------------------------------------------------

The simplest and most common mode.
The laser fires a single short pulse ($<50$\,ns), and a timer starts.
When the reflected echo arrives at the receiver, the timer stops.
The round-trip time $t$ gives the distance directly.

\begin{figure}[H]
  \centering
  \includegraphics[width=0.88\textwidth]{slide-11.jpg}
  \caption{Pulsed LiDAR timing diagram.  Top: the transmitted pulse train
    with period $T$ and amplitude $A_t$.  Bottom: the received echo train,
    delayed by the round-trip time $t$ and reduced in amplitude $A_R$.
    Right: physical setup --- the transmitter fires toward a surface at
    distance $R$; the receiver captures the return.}
\end{figure}

\textbf{Range from time of flight:}

The pulse travels to the target and back, covering a total distance of $2R$
at the speed of light $c$.  Solving for $R$:
\[
  R = \frac{1}{2}\,c\,t
\]
where $t$ is the measured round-trip time and $c \approx 3\times10^8$\,m/s.
The factor $\frac{1}{2}$ accounts for the \emph{two-way} trip --- the pulse
must travel to the target \emph{and} back.

\textbf{Distance resolution from sampling interval:}

In practice the timer is a digital clock with a fixed sampling interval
$\Delta t$ (the smallest time step it can resolve).  This limits how finely
the range can be measured.  The \emph{distance resolution} --- the smallest
difference in range that can be distinguished --- is:
\[
  \Delta R = \frac{1}{2}\,c\,\Delta t
\]
A shorter clock tick $\Delta t$ gives finer range resolution.  For example,
$\Delta t = 1$\,ns gives $\Delta R = 15$\,cm.

\begin{notebox}
  \textbf{What limits the maximum measurement distance?}

  The formula $R = \frac{1}{2}ct$ gives range from time alone, but the sensor
  can only detect a return if the echo is strong enough to register above the
  noise floor of the receiver.  The maximum range is therefore affected by:
  \begin{itemize}[leftmargin=1.2em,itemsep=2pt]
    \item \textbf{Laser power} --- more powerful pulse $\to$ stronger echo
      at long range.
    \item \textbf{Beam divergence} --- a wider beam spreads energy over a
      larger footprint at distance, reducing the power density hitting the
      target and the fraction reflected back.
    \item \textbf{Atmospheric transmission} --- air, water vapour, dust, and
      fog absorb and scatter the beam, attenuating it over distance.
    \item \textbf{Target reflectance} --- a dark or rough surface absorbs
      more of the pulse; a bright or smooth surface reflects more back.
    \item \textbf{Detector sensitivity} --- a more sensitive photodetector
      can register weaker returns.
  \end{itemize}
  All five factors feed into the \emph{LiDAR range equation} (the scattering
  model derived in Section~3).
\end{notebox}

% ------------------------------------------------------------------------------
\subsection*{2.2 \quad Phased LiDAR --- Continuous Wave with Phase Comparison}
% ------------------------------------------------------------------------------

Instead of firing a short pulse, phased (or phase-based) LiDAR emits a
\emph{continuous} laser beam whose \emph{intensity is modulated} sinusoidally
at a fixed frequency $f$.
The receiver compares the \emph{phase} of the returning wave with the phase of
the transmitted wave.
Because the light travels a round-trip distance $2R$ at speed $c$, the
returning wave is delayed by time $t = 2R/c$, which appears as a phase shift
$\phi$.

\begin{figure}[H]
  \centering
  \includegraphics[width=0.88\textwidth]{slide-12.jpg}
  \caption{Phased LiDAR signals.  Top: the transmitted sinusoidal wave with
    period $T = \lambda/c$ and amplitude $A_t$.  Bottom: the received wave,
    shifted in phase by $\phi$ relative to the transmitted wave.  The
    half-period $T/2$ corresponds to a round trip of one full wavelength
    $\lambda$.}
\end{figure}

\textbf{Phase difference to range:}

The wave completes one full cycle ($2\pi$ radians) in one period $T$.
The received signal is delayed by time $t$, which is a fraction $t/T$ of one
full cycle, corresponding to a phase shift $\phi$:
\[
  \frac{t}{T} = \frac{\phi}{2\pi}
  \quad\Longrightarrow\quad
  t = \frac{T\,\phi}{2\pi}
\]
Substituting into $R = \frac{1}{2}ct$ and using $T = 1/f = \lambda/c$:
\[
  R
  = \frac{1}{2}\,c\,\frac{T\,\phi}{2\pi}
  = \frac{c}{4\pi f}\,\phi
  = \frac{\lambda}{4\pi}\,\phi
\]
where $\lambda = c/f$ is the modulation wavelength (not the optical
wavelength of the laser itself, but the wavelength of the \emph{intensity
modulation}).

\begin{notebox}
  \textbf{Intuition for the phase formula.}

  Think of the modulated intensity as a ruler laid along the beam.  One full
  wavelength $\lambda$ of the modulation corresponds to a round-trip distance
  of $\lambda$ (i.e.\ a one-way range of $\lambda/2$).  Measuring the phase
  $\phi$ tells you \emph{where along that ruler} the target sits:
  \[
    R = \frac{\lambda}{4\pi}\,\phi
    \qquad
    \text{(range increases linearly with measured phase)}
  \]
  A full phase cycle $\phi = 2\pi$ gives $R = \lambda/2$ --- the maximum
  \emph{unambiguous} range.  Beyond this, the phase wraps around and the
  system cannot tell whether $\phi$ corresponds to one range or a range one
  modulation wavelength further away.  This is the fundamental range
  ambiguity of phased LiDAR.

  \begin{center}
    \begin{tabular}{@{}lll@{}}
      \toprule
      & \textbf{Pulsed LiDAR} & \textbf{Phased LiDAR} \\
      \midrule
      Range from    & Time of flight $t$              & Phase shift $\phi$ \\
      Formula       & $R = \tfrac{1}{2}ct$            & $R = \tfrac{\lambda}{4\pi}\phi$ \\
      Range limit   & Laser power / atmosphere        & Modulation wavelength ($\lambda/2$) \\
      Typical range & Long ($>100$\,m)                & Shorter ($\sim$1--100\,m) \\
      Precision     & Moderate                        & Higher \\
      \bottomrule
    \end{tabular}
  \end{center}

  \textbf{Why phased LiDAR is more precise.}
  Pulsed LiDAR measures a very short \emph{time interval} (nanoseconds);
  even a 1\,ns clock gives only $\Delta R = 15$\,cm resolution.
  Phased LiDAR instead measures a \emph{phase angle}, which electronic
  circuits (phase-locked loops, mixers) can resolve to fractions of a degree.
  Furthermore, the continuous wave \emph{averages over many cycles},
  suppressing noise the same way a longer integration time helps any
  measurement.  The result: sub-millimetre precision is achievable with
  phased LiDAR at short ranges, whereas pulsed systems rarely do better than
  a few centimetres.  The cost, as shown above, is the limited unambiguous
  range $\lambda/2$.
\end{notebox}

\begin{notebox}
  \textbf{Example calculation --- phased LiDAR range}

  Suppose the intensity modulation frequency is $f = 10$\,MHz.  Then:
  \[
    \lambda = \frac{c}{f} = \frac{3\times10^8}{10\times10^6} = 30\,\text{m}
    \qquad
    R_{\max} = \frac{\lambda}{2} = 15\,\text{m (unambiguous range)}
  \]
  If the receiver measures a phase shift of $\phi = 120^\circ = \tfrac{2\pi}{3}$\,rad:
  \[
    R = \frac{\lambda}{4\pi}\,\phi
      = \frac{30}{4\pi}\cdot\frac{2\pi}{3}
      = \frac{30}{6}
      = 5\,\text{m}
  \]
  The target is 5\,m away.

  \textbf{The ambiguity problem --- why the target must be closer than
  $\lambda/2$.}

  You are right to flag this.  The phase $\phi$ is always measured in the
  range $[0,\,2\pi)$ --- the receiver cannot count how many full cycles the
  wave has completed; it only sees \emph{where in the current cycle} the
  return sits.  This means that any target further than $\lambda/2$ produces
  a phase that \emph{wraps around} and looks identical to a closer target.

  \textit{Concrete example:} take the same $\lambda = 30$\,m system.
  \begin{itemize}[leftmargin=1.2em,itemsep=3pt]
    \item \textbf{Target A} at $R = 5$\,m:
      $\phi = \tfrac{4\pi \times 5}{30} = \tfrac{2\pi}{3} \approx 120^\circ$
      \quad$\to$ measured as $120^\circ$ \checkmark
    \item \textbf{Target B} at $R = 20$\,m:
      $\phi = \tfrac{4\pi \times 20}{30} = \tfrac{8\pi}{3} \approx 480^\circ$
      \quad$\to$ wraps: $480^\circ - 360^\circ = 120^\circ$
  \end{itemize}
  Target A at 5\,m and Target B at 20\,m give \textbf{exactly the same
  measured phase of $120^\circ$}.  The system cannot tell them apart.

  The wrap occurs at every multiple of $\lambda/2 = 15$\,m.  So 5\,m,
  20\,m, 35\,m, \ldots\ are all indistinguishable.

  \textbf{Design rule:} choose the modulation frequency $f$ so that
  $\lambda/2$ is \emph{larger} than the furthest target you ever expect:
  \[
    f < \frac{c}{2\,R_{\max,\,\text{expected}}}
  \]
  Lower frequency $\to$ longer $\lambda$ $\to$ larger unambiguous range,
  but also coarser phase ruler $\to$ worse precision.  This is the core
  trade-off of phased LiDAR.

  \textit{Note:} some systems use \textbf{two modulation frequencies}
  simultaneously.  The beat between them effectively creates a much longer
  ``virtual wavelength'', extending the unambiguous range while keeping
  fine precision from the higher frequency --- similar to a vernier scale.
\end{notebox}

% ------------------------------------------------------------------------------
\subsection*{2.3 \quad Ranging Precision}
% ------------------------------------------------------------------------------

\begin{figure}[H]
  \centering
  \includegraphics[width=0.82\textwidth]{slide-13.jpg}
  \caption{Ranging precision slide.  The two formulas give the standard
    deviation of the range estimate for pulsed and phased LiDAR respectively,
    showing how noise bandwidth and received power govern measurement quality.}
\end{figure}

\textbf{Ranging precision} is defined as the \emph{standard deviation}
$\sigma$ of the distance estimate in the presence of noise.
It is related to the \textbf{signal-to-noise ratio (SNR)} of the returned
laser signal: a stronger return relative to noise gives a more precise
distance estimate.

For the two LiDAR types, the precision scales as:

\[
  \sigma_{\text{pulse}}
  \;\propto\;
  \frac{c}{2}\,t_{\text{rise}}\,\sqrt{\frac{B_{\text{pulse}}}{P_{\text{peak}}}}
  \qquad\qquad
  \sigma_{\text{phase}}
  \;\propto\;
  \frac{\lambda}{4\pi}\,\sqrt{\frac{B_{\text{cw}}}{P_{\text{av}}}}
\]

\begin{notebox}
  \textbf{Variable definitions}

  \vspace{4pt}
  \begin{tabular}{@{}l p{4.2cm} p{7.2cm}@{}}
    \toprule
    \textbf{Symbol} & \textbf{Name} & \textbf{Meaning / effect on precision} \\
    \midrule
    $\sigma_{\text{pulse}}$ & Pulsed ranging std.\ dev.
      & Smaller = better precision \\[2pt]
    $\sigma_{\text{phase}}$ & Phased ranging std.\ dev.
      & Smaller = better precision \\[2pt]
    $c$ & Speed of light ($3\times10^8$\,m/s)
      & Fixed constant \\[2pt]
    $t_{\text{rise}}$ & Rise time of laser pulse
      & Shorter pulse edge $\to$ sharper timing $\to$ better precision \\[2pt]
    $\lambda$ & Modulation wavelength
      & Shorter $\lambda$ $\to$ finer phase ruler $\to$ better precision \\[2pt]
    $B_{\text{pulse}}$ & Noise bandwidth (pulsed)
      & Wider bandwidth captures more noise $\to$ worse precision \\[2pt]
    $B_{\text{cw}}$ & Noise bandwidth (phased/CW)
      & Same as above: wider = noisier \\[2pt]
    $P_{\text{peak}}$ & Peak power of received pulse
      & More power $\to$ stronger signal $\to$ better precision \\[2pt]
    $P_{\text{av}}$ & Average power of received CW signal
      & More power $\to$ better precision \\
    \bottomrule
  \end{tabular}
  \textbf{Understanding the key terms before the table}

  \textbf{Rise time $t_\text{rise}$ --- how sharp is the pulse edge?}
  The pulsed laser signal is not a sine wave --- it is a short rectangular
  burst (on/off), like a camera flash.  The rise time is simply how fast the
  leading edge switches from zero to full power.  Any sharp edge in time is
  equivalent to having high-frequency content in the signal (Fourier analysis:
  a perfectly sharp step would require infinitely high frequencies).  So:
  \[
    \text{shorter } t_\text{rise}
    \;\longleftrightarrow\;
    \text{steeper edge}
    \;\longleftrightarrow\;
    \text{higher frequency content in pulse}
  \]
  Because a blunt, slow-rising edge makes it hard to say exactly \emph{when}
  the pulse arrived, shorter $t_\text{rise}$ gives sharper timing and better
  range precision.

  \textbf{Noise bandwidth $B$ --- how much noise does the receiver let in?}
  Every receiver circuit responds to a range of frequencies --- its bandwidth.
  Noise (thermal, shot) is spread across \emph{all} frequencies, so a wider
  bandwidth lets more noise through alongside the signal, degrading precision.
  $B_\text{pulse}$ and $B_\text{cw}$ are the same concept but very different
  values in practice:
  \begin{itemize}[leftmargin=1.2em,itemsep=2pt]
    \item $B_\text{pulse}$ must be \emph{wide} --- a nanosecond pulse has
      frequency content up to $\sim$1\,GHz, so the pulsed receiver needs
      GHz-range bandwidth to capture the sharp edges faithfully.
    \item $B_\text{cw}$ can be \emph{narrow} --- only a small band around the
      modulation frequency (e.g.\ 10\,MHz) needs to pass; all other
      frequencies are noise and can be filtered out.
  \end{itemize}
  This is part of why phased LiDAR achieves better precision: its narrow
  $B_\text{cw}$ admits far less noise than the wide $B_\text{pulse}$ of a
  pulsed receiver, even at equal received power.

  \textbf{The ``ruler'' --- what physical distance corresponds to one unit of
  measurement?}
  Every measurement needs a scale.  In each LiDAR type, one unit of the
  measured quantity (time tick or phase radian) maps to a certain number of
  metres:
  \begin{itemize}[leftmargin=1.2em,itemsep=2pt]
    \item \textbf{Pulsed:} one clock tick $\Delta t = 1$\,ns
      $\to$ $\tfrac{1}{2}c\,\Delta t = 0.15$\,m.
      So 15\,cm is one graduation on the pulsed ruler.
      A shorter rise time shrinks this graduation.
    \item \textbf{Phased:} one radian of phase
      $\to$ $\tfrac{\lambda}{4\pi}$\,m.
      With $\lambda = 30$\,m that is $\sim$2.4\,m per radian --- sounds
      coarse, but electronics can resolve 0.001\,rad easily, giving
      $\sim$2.4\,mm precision.  A shorter modulation wavelength $\lambda$
      makes the ruler finer.
  \end{itemize}
  The $\sigma$ formulas simply say: precision = (ruler graduation)
  $\times$ (noise-to-signal ratio).


  \vspace{6pt}
  The key insight is the \textbf{square-root SNR dependence}:
  precision scales as $\sqrt{B/P}$.  To halve $\sigma$ (double the
  precision), you need to \emph{quadruple} the received power, or quarter the
  noise bandwidth --- improvements are expensive because of the square root.

  Notice also the \textbf{structural similarity} between the two formulas:
  \begin{itemize}[leftmargin=1.2em,itemsep=2pt]
    \item Pulsed: the ``ruler'' is set by $\tfrac{c}{2}\,t_{\text{rise}}$
      --- essentially the spatial length of the pulse edge in metres.
      The shorter and sharper the pulse edge, the finer the timing and the
      better the precision.
    \item Phased: the ``ruler'' is set by $\tfrac{\lambda}{4\pi}$
      --- one radian of phase in metres.  A shorter modulation wavelength
      gives a finer ruler, hence better precision.  This is why phased LiDAR
      typically outperforms pulsed LiDAR in precision at short ranges.
  \end{itemize}
\end{notebox}

% ── Key takeaways ─────────────────────────────────────────────────────────────
\begin{keybox}
  \textbf{Section 2 --- Key Takeaways}
  \begin{itemize}[leftmargin=1.2em,itemsep=2pt]
    \item \textbf{Pulsed LiDAR} fires a short burst and measures round-trip
      time: $R = \tfrac{1}{2}ct$.  Distance resolution is set by the clock
      interval: $\Delta R = \tfrac{1}{2}c\,\Delta t$.
    \item \textbf{Phased LiDAR} modulates a continuous beam and measures the
      phase shift of the return: $R = \tfrac{\lambda}{4\pi}\phi$.
      Higher precision but limited unambiguous range ($\leq \lambda/2$).
    \item \textbf{Ranging precision} $\sigma \propto \sqrt{B/P}$: halving
      $\sigma$ requires quadrupling received power.  Precision improves with
      sharper pulse edges (pulsed) or shorter modulation wavelength (phased).
    \item The \textbf{maximum range} is governed by physics beyond the timing
      formula: laser power, beam divergence, atmosphere, target reflectance,
      and detector sensitivity all play a role (covered fully in Section~3).
  \end{itemize}
\end{keybox}

% ==============================================================================
\newpage
\section*{Section 3 \quad LiDAR Scattering \& Waveform Model}
% ==============================================================================

We now know \emph{how} LiDAR measures time or phase to get distance.
But there is a deeper question: \emph{how much} of the emitted laser power
actually makes it back to the sensor?
This matters because if the returned power is too weak the measurement fails
entirely --- range is not infinite.
The \textbf{scattering model} traces the laser power step by step from
transmitter to target to receiver, giving a formula for the received power
$P_R$.
The \textbf{waveform model} then explains what happens when a single pulse
hits a complex target such as a tree, producing multiple overlapping returns.

% ------------------------------------------------------------------------------
\subsection*{3.1 \quad The Four-Step Power Chain}
% ------------------------------------------------------------------------------

\begin{figure}[H]
  \centering
  \includegraphics[width=0.82\textwidth]{slide-14.jpg}
  \caption{LiDAR scattering geometry.  The laser illuminates a footprint of
    area $A_\text{laser}$ on the ground at range $R$.  The target scatters
    power in all directions; only the fraction caught by the receiver aperture
    $A_r = \pi D_r^2/4$ returns to the sensor.  The four numbered steps track
    the power budget from transmitter to receiver.}
\end{figure}

It helps to follow the power through four physical steps:

\begin{enumerate}[leftmargin=1.6em,itemsep=4pt]
  \item \textbf{Laser emits power $P_t$.}
    The transmitter sends out a pulse with total power $P_t$.
    This power spreads over a cone as it travels downward.

  \item \textbf{Power incident on the target.}
    At range $R$, the beam illuminates a circular footprint of area
    $A_\text{laser}$ on the ground.
    $A_\text{laser}$ is the \emph{total} area lit up by the beam --- it can
    be large (a wide footprint at long range hits many things: ground, rocks,
    vegetation).
    The target itself only occupies part of that footprint; its
    \emph{effective scattering area} $A_S \leq A_\text{laser}$ is how much
    of the footprint it fills.
    A small rock sitting inside a large footprint has $A_S \ll A_\text{laser}$
    --- only a tiny fraction of the emitted power actually hits it.
    The ratio $A_S / A_\text{laser}$ is therefore the fraction of $P_t$ that
    reaches the target:
    \[
      P_\text{incident} = P_t \cdot \frac{A_S}{A_\text{laser}}
    \]

  \item \textbf{Target scatters power.}
    The \emph{reflectance} $\rho$ ($0 \leq \rho \leq 1$) is primarily a
    \textbf{material property} --- it describes what fraction of the incident
    light the surface reflects rather than absorbs.
    Dark asphalt has $\rho \approx 0.05$; white paint has $\rho \approx 0.9$.
    The shape and roughness of the surface determine the \emph{direction} the
    light is scattered (captured by $\Omega$), but $\rho$ itself is the
    material's reflective fraction --- the two effects are kept separate in
    the model.
    The scattered power is:
    \[
      P_S = P_t \cdot \frac{A_S}{A_\text{laser}} \cdot \rho
    \]

  \item \textbf{Receiver captures a fraction of the scattered power.}
    $\Omega$ is the \emph{full cone of directions} the target scatters light
    into --- it has nothing to do with where the sensor is.  Three cases
    illustrate the range:
    \begin{itemize}[leftmargin=1.2em,itemsep=2pt]
      \item \textbf{Mirror / retroreflector} --- reflects nearly all light
        back in one narrow direction.  $\Omega$ is tiny; most scattered power
        lands on the sensor $\to$ very strong return.  This is why road signs
        and survey targets coated with retroreflective material give extremely
        bright LiDAR returns.
      \item \textbf{Rough / diffuse surface} (soil, bark, asphalt) ---
        scatters light roughly equally in all directions of the upper
        hemisphere.  $\Omega \approx 2\pi$ steradians.  The sensor intercepts
        only a tiny slice $\to$ weak return.
      \item \textbf{Most real surfaces} sit somewhere between: a preferred
        specular peak plus a broad diffuse component.
    \end{itemize}
    The formula captures this directly: received fraction $= A_r /
    (\Omega R^2)$.  Smaller $\Omega$ (more directed scatter) $\to$ larger
    fraction back to the sensor.

    The scattered power spreads over solid angle $\Omega$ at range $R$,
    covering an area of $\Omega R^2$.
    The receiver lens has diameter $D_r$, so its collecting area is:
    \[
      A_r = \frac{\pi D_r^2}{4}
      \quad\text{(area of a circle with diameter } D_r\text{)}
    \]
    This is simply $\pi r^2$ with $r = D_r/2$ --- a bigger lens intercepts
    more of the scattered sphere.
    The fraction captured is $A_r / (\Omega R^2)$.

    The two efficiency factors $\eta_\text{atm}$ and $\eta_\text{sys}$ cover
    losses on \textbf{both the outgoing and return paths}, not just the
    return.  They appear once at the end because $P_t$ is defined as the
    power at the laser output; all subsequent losses --- atmosphere on the
    way \emph{to} the target, atmosphere on the way \emph{back}, plus sensor
    optics and electronics --- are multiplied in together as a combined
    round-trip efficiency.
    (If you wanted to split them: $\eta_\text{atm} = \eta_\text{atm,go}
    \times \eta_\text{atm,return}$, but in practice a single factor is
    measured and used.)
    The received power is:
    \[
      P_R = P_t \cdot \frac{A_S}{A_\text{laser}} \cdot \rho
            \cdot \frac{1}{\Omega R^2} \cdot \frac{\pi D_r^2}{4}
            \cdot \eta_\text{atm}\,\eta_\text{sys}
    \]
\end{enumerate}

% ------------------------------------------------------------------------------
\subsection*{3.2 \quad Deriving the Full Scattering Model}
% ------------------------------------------------------------------------------

\begin{figure}[H]
  \centering
  \includegraphics[width=0.82\textwidth]{slide-15.jpg}
  \caption{LiDAR scattering model with labelled geometry.  $\beta_t$ is the
    beam divergence half-angle; the illuminated footprint grows as $R^2$.
    The scatterer disperses power uniformly into solid angle $\Omega$.}
\end{figure}

Two assumptions are made to make the model tractable:
\begin{enumerate}[leftmargin=1.6em,itemsep=2pt]
  \item The emitted laser power is \textbf{uniformly distributed} within the
    footprint (flat-top beam profile).
  \item The scatterer \textbf{uniformly disperses} the incident power into
    a conical solid angle $\Omega$ --- a so-called \textbf{Lambertian}
    surface model.
\end{enumerate}

\begin{warnbox}
  \textbf{Assumption 2 is a simplification --- real surfaces do not scatter
  uniformly.}

  The Lambertian assumption says power is spread \emph{evenly} inside
  $\Omega$.  Real surfaces do not do this.  Ridges, bumps, and surface
  texture mean some scatter directions get more energy than others.  Most
  surfaces also have a \textbf{specular component} --- a stronger peak of
  reflection in the mirror direction on top of the broad diffuse scatter.
  The angle the beam hits the surface at further changes the pattern.

  The full description of how a real surface scatters is called the
  \textbf{BRDF} (Bidirectional Reflectance Distribution Function), which
  gives a different value for every incoming/outgoing angle combination.

  The model uses Lambertian scattering for three practical reasons:
  \begin{itemize}[leftmargin=1.2em,itemsep=2pt]
    \item The BRDF of every target in the scene is unknown in advance.
    \item Uniform $\Omega$ makes the maths tractable and gives one
      closed-form formula for $P_R$.
    \item For most rough natural surfaces (soil, vegetation, asphalt) the
      non-uniformities average out over many measurements, so the
      Lambertian model is a reasonable mean estimate.
  \end{itemize}
  $\Omega$ therefore acts as a single model parameter that lumps all unknown
  directional scattering behaviour into one number.  The formula gives a
  useful engineering estimate, not an exact value.
\end{warnbox}

\textbf{Step 1 --- How large is the illuminated footprint at range $R$?}

The laser beam is not a perfectly parallel ray --- it spreads slightly as
it travels, forming a cone.
The \textbf{half-angle} $\varphi$ is the angle between the centre axis of
the beam and its outer edge, i.e.\ half the total opening angle of the cone.
Think of a flashlight: a wide half-angle gives a broad circle of light on
the wall; a narrow half-angle gives a tight spot.

At range $R$, the beam has spread to a circular footprint of radius
$r = R\tan\varphi \approx R\beta_t/2$, so its area is:
\[
  A_\text{laser}
  = \pi r^2
  = \frac{\pi R^2 \tan^2\!\varphi}{4}
  = \frac{\pi R^2 \beta_t^2}{4}
\]
\textit{What this tells you:} $A_\text{laser}$ is what we called the
footprint area back in Step~2 of the four-step chain.  Now we have an
explicit formula for it.  Because $A_\text{laser} \propto R^2$, doubling
the range quadruples the footprint --- the same total power $P_t$ is
spread over four times the area, so the power \emph{density} on the ground
drops to one quarter.  A narrow beam (small $\beta_t$) keeps the footprint
small and the power density high even at long range.

\begin{warnbox}
The beam itself is \textbf{3D} --- a cone expanding in all directions as it
travels.  The footprint it makes on the ground is \textbf{2D} --- a flat
circle.  The formula $\pi R^2\beta_t^2/4$ assumes the \emph{same}
half-angle $\beta_t$ in every direction around the beam axis, i.e.\ the
cone is perfectly circular in cross-section.  If the divergence were
different in the two perpendicular directions ($\beta_x \neq \beta_y$),
the footprint would instead be an \textbf{ellipse} with area
$\pi R^2 \beta_x \beta_y / 4$.  Most LiDAR systems are designed for a
circular beam, so $\beta_x = \beta_y = \beta_t$ is a reasonable
assumption.
\end{warnbox}

\textbf{Step 2 --- How much power reaches the sensor after scattering?}

We now substitute the explicit $A_\text{laser}$ from Step~1 into the
scattered-power formula from Step~3 of the four-step chain
($P_S = P_t \cdot \tfrac{A_S}{A_\text{laser}} \cdot \rho$):
\[
  P_S
  = P_t \cdot \frac{A_S}{\dfrac{\pi R^2 \beta_t^2}{4}} \cdot \rho
  = \frac{4\,P_t}{\pi R^2 \beta_t^2}\,\rho\,A_S
\]
\textit{What this tells you:} scattered power already falls off as
$1/R^2$ --- the bigger footprint dilutes the power hitting the target.
This uses the Lambertian model from assumption~2: the target disperses
the incident power uniformly across $\Omega$, so the fraction that
eventually returns depends purely on $\rho$, $A_S$, and how large
$A_\text{laser}$ is at that range.

\textbf{Step 3 --- What power finally arrives at the receiver?}

We now take the full expression from Step~4 of the four-step chain and
substitute $A_\text{laser}$ everywhere:
\[
  \boxed{
  P_R
  = \frac{4\,P_t}{\pi R^2 \beta_t^2}
    \cdot \rho A_S
    \cdot \frac{1}{\Omega R^2}
    \cdot \frac{\pi D_r^2}{4}
    \cdot \eta_\text{atm}\,\eta_\text{sys}
  }
\]
\textit{What this tells you:} this is the complete \textbf{LiDAR range
equation}.  The two $1/R^2$ factors combine to give a $1/R^4$ penalty ---
doubling the range reduces returned power by $16\times$.  The formula
also makes clear which design choices help: a narrower beam ($\beta_t$
smaller), a larger receiver lens ($D_r$ larger), a more reflective or
larger target ($\rho A_S$ larger), and low atmospheric loss ($\eta_\text{atm}$
close to~1) all directly increase $P_R$.  This is exactly what limits the
maximum range mentioned back in Section~2 (the five limiting factors in the
pulsed LiDAR notebox) --- they all appear here explicitly.

\begin{notebox}
  \textbf{Variable definitions}

  \vspace{4pt}
  \begin{tabular}{@{}l p{3.8cm} p{7.2cm}@{}}
    \toprule
    \textbf{Symbol} & \textbf{Name} & \textbf{Meaning} \\
    \midrule
    $P_t$           & Transmitted power
      & Total laser pulse power emitted \\[2pt]
    $R$             & Range
      & Distance from sensor to target \\[2pt]
    $\beta_t$       & Beam divergence
      & Half-angle of the emitted cone; wider = larger footprint \\[2pt]
    $A_\text{laser}$& Footprint area
      & Illuminated area at range $R$; grows as $R^2$ \\[2pt]
    $A_S$           & Scatterer area
      & Effective reflecting area of the target \\[2pt]
    $\rho$          & Reflectance
      & Fraction of incident power reflected ($0$--$1$); dark surfaces $\approx 0$, bright $\approx 1$ \\[2pt]
    $\Omega$        & Solid angle
      & Cone into which the target disperses the power \\[2pt]
    $D_r$           & Receiver diameter
      & Aperture of the sensor lens; larger = captures more \\[2pt]
    $\eta_\text{atm}$ & Atmospheric efficiency
      & Loss due to absorption and scattering in air ($\leq 1$) \\[2pt]
    $\eta_\text{sys}$ & System efficiency
      & Loss in the sensor's own optics and electronics ($\leq 1$) \\
    \bottomrule
  \end{tabular}

  \vspace{6pt}
  \textbf{Key observation: the $R^4$ penalty.}
  Notice that $P_R \propto 1/R^4$ --- the received power drops as the
  \emph{fourth power} of range.  The $1/R^2$ comes from the footprint
  growing (less power hits the target), and another $1/R^2$ comes from
  the scattered power spreading over the hemisphere before the receiver
  catches its small slice.  Doubling the range reduces returned power by
  a factor of $16$, which is why long-range LiDAR requires powerful lasers
  and sensitive detectors.
\end{notebox}

% ------------------------------------------------------------------------------
\subsection*{3.3 \quad Multiple Returns --- The ``But\ldots'' Slide}
% ------------------------------------------------------------------------------

\begin{figure}[H]
  \centering
  \includegraphics[width=0.82\textwidth]{slide-16.jpg}
  \caption{A single LiDAR pulse hitting a tree.  The beam illuminates a
    footprint that includes both foliage and ground.  The pulse partially
    reflects off the tree canopy (first return), travels further to lower
    branches (second return), and finally reflects off the ground (ground
    return).  The receiver records a time-series of overlapping pulses rather
    than a single clean echo.}
\end{figure}

The scattering model above assumed a single flat target at range $R$.
In practice, a single pulse often illuminates a \emph{complex} scene
--- a tree, a building with windows, or uneven terrain.
Different parts of the scene are at different ranges, so the pulse
reflects back at different times, and the receiver records a stretched,
multi-peaked waveform instead of a single clean echo.

As shown in the figure, one pulse hitting a tree produces three peaks
in the received waveform:
\begin{itemize}[leftmargin=1.4em,itemsep=2pt]
  \item A \textbf{canopy return} --- the first reflection from the top
    of the tree.  Arrives earliest (shortest range).
  \item A \textbf{low canopy / branch return} --- a weaker echo from
    lower branches that the pulse partially penetrated.
  \item A \textbf{ground return} --- the remaining pulse energy that
    passed through all gaps in the foliage and reflected off the ground.
    Arrives last (longest range).
\end{itemize}

Each peak in the waveform has a \textbf{height} (amplitude) and a
\textbf{width}.  Here is what each tells you:

\textbf{Peak height --- how much power came back from that surface.}
A taller peak means more power returned.  But the height is governed
by the full scattering formula $P_R \propto \rho\, A_S / R^4$, so a
shorter peak could mean:
\begin{itemize}[leftmargin=1.2em,itemsep=2pt]
  \item A weakly-reflecting surface (small $\rho$), \emph{or}
  \item A small surface area (small $A_S$), \emph{or}
  \item Simply a surface \textbf{farther away} (large $R$, $R^4$
    scaling).
\end{itemize}
In the tree example: the ground return peak is typically the weakest
of the three, but this does \textbf{not} mean the ground is a poor
reflector.  The ground is simply the furthest surface --- the pulse
has already lost energy passing through canopy and branches, and the
remaining path is longest.  You cannot conclude anything about
reflectance from height alone without first accounting for range.

\textbf{Peak width --- what the surface looks like.}
The emitted pulse has a finite duration, so every return has a minimum
width.  Beyond that, the echo is broadened by:
\begin{itemize}[leftmargin=1.2em,itemsep=2pt]
  \item \textbf{Surface roughness} --- height variations within the
    footprint mean different parts reflect at slightly different ranges,
    smearing the echo in time.
  \item \textbf{Surface tilt} --- a tilted surface has one edge closer
    than the other, broadening the return.
  \item \textbf{Depth extent} --- a thick layer (dense foliage) reflects
    from both front and back, producing a wide merged peak.
\end{itemize}
In the tree example: the canopy return is often broad (the canopy has
depth and rough texture); the ground return is typically narrow and
sharp (flat, smooth ground reflects cleanly).

\textbf{You need timing, height, and width together.}
The peak arrival time gives range; height corrected for that range gives
reflectance; width gives surface character.  None of the three is
meaningful in isolation --- a low, broad peak near the sensor could
be a rough nearby surface, while an equally low but narrow peak far
away is likely a smooth flat surface simply attenuated by $R^4$.
This is why full-waveform LiDAR records the \emph{entire} time-series
rather than just the first return time.

% ------------------------------------------------------------------------------
\subsection*{3.4 \quad LiDAR Waveform Model}
% ------------------------------------------------------------------------------

\begin{figure}[H]
  \centering
  \includegraphics[width=0.82\textwidth]{slide-17.jpg}
  \caption{LiDAR waveform model.  The emitted pulse is approximately
    Gaussian in time.  Convolution with the scattering transfer function
    (one Gaussian per scatterer) produces the full received waveform, which
    is a sum of Gaussians --- one per distinct reflecting surface.}
\end{figure}

This model applies to \textbf{pulsed} LiDAR --- the sensor fires discrete
pulses and records the returned power as a function of time.
Phased LiDAR instead compares the phase of a continuous wave and does
not produce a time-series waveform of this kind.

\textbf{Why the emitted pulse is Gaussian in time.}
When the laser fires, the amplitude does not switch on and off
instantaneously --- it ramps up, reaches a peak, and ramps back down.
The rise and fall are smooth, making the pulse shape approximately
\textbf{Gaussian}: a bell curve in time with the peak at the centre of
the pulse duration.
There is no perfectly sharp pulse in practice; finite rise time (recall
Section~2) is exactly this gradual ramp.

\textbf{Why each surface return is also Gaussian.}
When this Gaussian pulse hits a single flat surface at range $R_i$, the
surface reflects a delayed, scaled copy of the incoming pulse back toward
the sensor.
The copy arrives delayed by $2R_i/c$ and scaled by the surface's
reflectance and geometry --- but its \emph{shape} is unchanged.
Since the incoming pulse is Gaussian, the reflected copy is also Gaussian.

\textbf{Convolution --- why we use that word.}
Mathematically, the received waveform is the \emph{convolution} of the
emitted pulse with the \emph{scattering transfer function}.
The scattering transfer function describes the scene: it has a spike at
each time $\mu_i = 2R_i/c$ with height $A_i$ for each reflecting surface.
Convolving the emitted Gaussian with each spike shifts and scales a copy
of the Gaussian to that time --- one copy per surface.
The key property that makes the model clean: \textbf{Gaussian $\ast$
Gaussian = Gaussian}, so each return stays Gaussian regardless of how
many surfaces there are.

\textbf{The full waveform is the sum of all these Gaussians.}
With $N$ reflecting surfaces, the receiver sees all their reflected copies
arriving on top of each other.
The total received power at any time $t$ is just their sum:

\[
  P_R(t)
  = \sum_{i=1}^{N_\text{scatterers}}
    A_i \, \exp\!\left(-\frac{(t - \mu_i)^2}{2\sigma_i^2}\right)
\]
Each term in the sum is one Gaussian peak: one surface, one range, one
echo.  The full waveform is the overlay of all of them.

\begin{notebox}
  \textbf{Variable definitions and what each tells you}

  \vspace{4pt}
  \begin{tabular}{@{}l p{3.5cm} p{7.5cm}@{}}
    \toprule
    \textbf{Symbol} & \textbf{Name} & \textbf{Meaning} \\
    \midrule
    $P_R(t)$   & Received power waveform
      & Full time-series of returned laser power at the receiver \\[2pt]
    $N_\text{scatterers}$ & Number of scatterers
      & How many distinct reflecting surfaces the pulse hit \\[2pt]
    $A_i$      & Amplitude of return $i$
      & How strongly surface $i$ reflected; related to $\rho$ and $A_S$ \\[2pt]
    $\mu_i$    & Mean (centre time) of return $i$
      & Arrival time of the $i$-th echo $\to$ range via $R_i = \tfrac{1}{2}c\mu_i$ \\[2pt]
    $\sigma_i$ & Width of return $i$
      & How spread-out the echo is; affected by pulse width, surface
        roughness, and slope of the target \\
    \bottomrule
  \end{tabular}

  \vspace{6pt}
  By fitting this model to the measured waveform you recover $A_i$,
  $\mu_i$, $\sigma_i$ for every peak simultaneously --- range,
  reflectance, and surface character from a single shot.
  See Section~3.3 for what each parameter physically means in the
  context of the canopy example.
\end{notebox}

% ── Key takeaways ─────────────────────────────────────────────────────────────
\begin{keybox}
  \textbf{Section 3 --- Key Takeaways}
  \begin{itemize}[leftmargin=1.2em,itemsep=2pt]
    \item Received power follows a \textbf{four-step chain}: emit
      $\to$ illuminate footprint $\to$ scatter at target $\to$ receiver
      catches a fraction.
    \item The full formula is
      $P_R \propto \rho A_S / (\beta_t^2 R^4)$ ---
      received power drops as $R^4$.
      Doubling range means 16$\times$ less power returned.
    \item \textbf{Footprint grows as $R^2$}: beam divergence $\beta_t$
      determines how quickly the illuminated area expands with range.
    \item A single pulse hitting a tree (or any layered scene) produces
      \textbf{multiple returns} at different times --- canopy, branches,
      ground.  These appear as separate Gaussian peaks in the waveform.
    \item The received waveform is modelled as a
      \textbf{sum of Gaussians}:
      $P_R(t) = \sum_i A_i \exp\!\bigl(-(t-\mu_i)^2 / 2\sigma_i^2\bigr)$.
      Fitting this model extracts range, reflectance, and surface roughness
      for each scatterer simultaneously.
  \end{itemize}
\end{keybox}

% ==============================================================================
\newpage
\section*{Section 4 \quad Point Cloud Data \& Real Hardware}
% ==============================================================================

The previous sections explained how LiDAR measures range and how the
returned power is modelled.
The output of a LiDAR sensor is not a single distance --- it is a
\textbf{point cloud}: a collection of thousands or millions of 3D points,
each carrying its own position, timestamp, and intensity value.
This section covers what a point cloud actually contains, how it is stored,
and what a real-world rotating LiDAR system looks like in practice.

% ------------------------------------------------------------------------------
\subsection*{4.1 \quad What a Point Cloud Contains}
% ------------------------------------------------------------------------------

A \textbf{point cloud} is simply a large collection of individual 3D points,
where each point represents one surface location that was measured by the
sensor.
Think of it as a scatter plot in 3D space: the sensor fires thousands of laser
pulses per second, each returning a range and direction, and every single
measurement becomes one dot in the cloud.
Enough dots together and the shape of the whole environment emerges ---
walls, ground, trees, vehicles --- even though there is no continuous surface,
just a dense set of discrete points.

The point cloud is \textbf{not stored on the sensor itself} --- the sensor
has no meaningful on-board memory.
Instead it streams the raw measurements in real time to a connected computer
(in the VLP-16 case, via Ethernet UDP packets, as covered in Section~4.2).
The computer then assembles the stream into a point cloud file on disk in a
standard format such as LAS or LAZ.

\begin{figure}[H]
  \centering
  \includegraphics[width=0.82\textwidth]{slide-18.jpg}
  \caption{Point cloud data structure.  Each individual point in the cloud
    carries several pieces of information beyond just its 3D position.}
\end{figure}

Each individual point in a point cloud contains:

\begin{itemize}[leftmargin=1.4em,itemsep=3pt]
  \item \textbf{3D coordinates} --- either Cartesian $(x, y, z)$ or polar
    $(R, \alpha, \theta)$ (range, azimuth, elevation).
    In a rotating LiDAR the sensor natively measures polar coordinates;
    these are then converted to Cartesian using the rotary encoder angles
    recorded at the moment of each shot.

  \item \textbf{Timestamp} --- the exact time the pulse was fired.
    This is critical for moving platforms (robots, cars, aircraft): without
    a timestamp, points recorded at different moments during motion cannot
    be correctly placed in a global coordinate frame.

  \item \textbf{Return signal strength (intensity)} --- the amplitude $A_i$
    of the returned pulse, as discussed in Section~3.
    This is the ``colour'' value of the point: brighter = more reflective
    surface.
    See the coloured-intensity image below (slide~19).

  \item \textbf{Classification results} --- some systems attach a label to
    each point indicating what type of object it likely belongs to
    (e.g.\ ground, vegetation, building, vehicle).
    This is the output of a post-processing step, not a raw sensor
    measurement.
\end{itemize}
\begin{notebox}
  \textbf{Point cloud file formats}

  Point cloud data can be stored in several standard formats:
  \begin{itemize}[leftmargin=1.2em,itemsep=2pt]
    \item \textbf{LAS} --- the most widely used industry standard for storing
      LiDAR point clouds.  Stores 3D coordinates, intensity, timestamp,
      classification, and other attributes in a compact binary format.
    \item \textbf{LAZ} --- a lossless compressed version of LAS.  Typically
      5--10$\times$ smaller file size with no data loss; the standard choice
      for archiving and transfer.
    \item \textbf{E57} --- a vendor-neutral format for 3D imaging data
      including point clouds.  Designed to be hardware-agnostic, so data
      from any scanner can be stored without proprietary dependencies.
  \end{itemize}
  Note: some sensors use \emph{proprietary} formats (e.g.\ Velodyne's PCAP)
  that require knowledge of the sensor's own geometry to convert to a
  standard format.
\end{notebox}

\begin{figure}[H]
  \centering
  \includegraphics[width=0.80\textwidth]{slide-19.jpg}
  \caption{Aerial LiDAR point cloud coloured by return intensity.  Bright
    points = highly reflective surfaces (road markings, vehicle roofs).
    Dark points = absorbing surfaces (vegetation, dark asphalt).  The car
    on the forest track is clearly visible due to its high reflectivity
    relative to the surrounding canopy.  Note: this is \emph{not} real
    colour --- it is a false-colour map of a single intensity number per
    point.}
\end{figure}

\begin{notebox}
  \textbf{The image looks coloured --- but there is no real colour in LiDAR}

  This is a common source of confusion.
  A standard LiDAR sensor such as the VLP-16 uses a \textbf{single laser
  wavelength} (903\,nm near-infrared).
  It records \emph{one number} per point --- the return intensity $A_i$
  (how much laser power came back).
  There is no red, green, or blue channel; no colour information whatsoever.

  What you see in the image is a \textbf{false-colour map}: the software
  takes the intensity number and maps it to a colour scale
  (e.g.\ high intensity $\to$ bright/white, low intensity $\to$ dark/black).
  It is exactly like colouring a terrain map by elevation --- the colours
  are chosen to make differences visible, they do not represent physical
  colours of the surfaces.

  The differences in brightness come entirely from \textbf{how reflective
  each surface is at 903\,nm near-infrared}:
  \begin{itemize}[leftmargin=1.2em,itemsep=2pt]
    \item Road markings, metal surfaces, vehicle roofs $\to$ strong
      reflectors $\to$ bright
    \item Dark asphalt, vegetation $\to$ absorb near-IR $\to$ dark
    \item The car stands out brightly against the dark canopy for exactly
      this reason --- not because it has a different colour, but because
      its surface reflects far more 903\,nm laser power
  \end{itemize}

  The multi-wavelength LiDAR mentioned in Section~1.3 is a separate,
  more advanced technique where two wavelengths are fired simultaneously
  and the \emph{ratio} of the two intensity values is used to classify
  materials.  But that is not what is shown here --- this image is purely
  single-wavelength intensity, false-coloured for visualisation.
\end{notebox}


% ------------------------------------------------------------------------------
\subsection*{4.2 \quad Real Hardware: Velodyne VLP-16}
% ------------------------------------------------------------------------------

\begin{figure}[H]
  \centering
  \includegraphics[width=0.88\textwidth]{slide-20.jpg}
  \caption{Velodyne VLP-16 ``Puck'' system setup.  The sensor connects to a
    laptop via Ethernet (UDP packets).  An optional INS/GPS unit provides
    georeferencing.  The power supply runs from 9--18\,V.}
\end{figure}

The \textbf{Velodyne VLP-16} (nicknamed the ``Puck'') is a compact
rotating LiDAR sensor widely used in robotics and autonomous vehicles.
Its key specifications:

\begin{itemize}[leftmargin=1.4em,itemsep=3pt]
  \item \textbf{16 laser/detector pairs} stacked at 16 different fixed
    \textbf{elevation angles}, ranging from $-15^\circ$ to $+15^\circ$
    in 2$^\circ$ steps.
    As the head spins 360$^\circ$ horizontally, all 16 lasers fire at
    every azimuth step --- not one per rotation, but all 16 \emph{on
    every shot}.
    One full rotation therefore produces \textbf{16 concentric rings} of
    points, one ring per elevation angle, stacked vertically on top of
    each other.
    More lasers = more vertical rings = denser coverage of the scene
    height.
    The head completes roughly 10--20 full rotations per second.

  \item \textbf{Fire rate of 18.08\,kHz} --- each laser fires 18,080 times
    per second.  With 16 lasers, this gives up to
    $16 \times 18{,}080 \approx \mathbf{300{,}000}$ \textbf{data points
    per second}.

  \item \textbf{Data packets every 1.3\,ms} --- by default, data
    accumulates over 24 firing cycles before being sent as a UDP packet
    over Ethernet.  This means the host computer receives a burst of
    points every 1.3\,ms.

  \item \textbf{Interface: Ethernet (UDP)} --- the sensor streams data
    directly over a standard network cable.  No special hardware is
    needed on the host side.

  \item \textbf{Storage: PCAP format} --- the raw UDP stream is typically
    captured in PCAP (packet capture) format.  Converting to LAS or
    another standard format requires precise knowledge of the sensor's
    own mounting position, orientation, and timing, since the PCAP
    contains raw firing angles and distances --- not fully-formed 3D
    coordinates.

  \item \textbf{Optional INS/GPS} --- an Inertial Navigation System
    (INS) measures the sensor's acceleration and rotation continuously,
    tracking its position and orientation at all times.
    GPS adds absolute position in the world.
    Together they answer: \emph{where was the sensor, and which way was
    it pointing, at the exact moment each pulse fired?}

    This is \textbf{optional for a static sensor} --- if the sensor sits
    on a tripod at a known fixed location, its position and orientation
    are already known and no INS/GPS is needed.

    It becomes \textbf{essential for a moving platform} (robot, car,
    aircraft): without knowing the sensor's exact position at each
    firing moment, points recorded at different instants during motion
    get placed in wrong locations and the resulting map smears or
    distorts.
    The timestamp on each measurement packet is what links a LiDAR point
    to the correct INS/GPS position at that exact moment.
\end{itemize}

% ------------------------------------------------------------------------------
\subsection*{4.3 \quad VLP-16 Data Packet Structure}
% ------------------------------------------------------------------------------

\begin{figure}[H]
  \centering
  \includegraphics[width=0.90\textwidth]{slide-21.jpg}
  \caption{VLP-16 UDP data packet structure.  Each packet is 1248 bytes
    and contains two firing sequences, each with 12 data blocks.  Every
    data block stores one azimuth angle and 32 channel measurements
    (distance + reflectivity per channel).}
\end{figure}

Each UDP packet sent by the VLP-16 has a fixed size of \textbf{1248 bytes},
structured as follows:

\begin{itemize}[leftmargin=1.4em,itemsep=3pt]
  \item \textbf{42-byte header} --- network and protocol overhead (IP/UDP
    headers).

  \item \textbf{1200 bytes of data} --- $12 \times 2$ bytes (flag) $+$
    $2$ bytes (azimuth) $+$ $32 \times (2$ bytes distance $+$ 1 byte
    reflectivity$)$ per data block, across 12 blocks per firing sequence
    and 2 firing sequences per packet.
    Each data block contains one azimuth reading and 32 channel
    measurements (distance and reflectivity for each of the 16 lasers,
    two returns per laser).

  \item \textbf{4-byte timestamp} --- GPS-synchronised time of the packet.

  \item \textbf{2-byte factory field} --- sensor model and return mode
    identifier.
\end{itemize}

\begin{notebox}
  \textbf{Why does the packet structure matter?}

  When you capture raw PCAP data from the VLP-16, you receive a stream
  of these 1248-byte packets.  To reconstruct a 3D point cloud you must:
  \begin{enumerate}[leftmargin=1.4em,itemsep=2pt]
    \item Parse each packet to extract the azimuth angle and the 32
      distance/reflectivity values per block.
    \item Know the \emph{vertical angles} of each of the 16 laser beams
      (factory-calibrated, stored in the sensor's own memory).
    \item Convert each (range, azimuth, elevation) polar measurement to
      Cartesian $(x, y, z)$.
    \item Apply the timestamp and any platform motion correction
      (from INS/GPS) to place points in a global frame.
  \end{enumerate}
  This is why conversion from PCAP to LAS requires ``precise knowledge
  about the sensor's own position'' as stated on the slide --- the raw
  packet contains angles and distances, not world coordinates.
\end{notebox}

% ── Key takeaways ─────────────────────────────────────────────────────────────
\begin{keybox}
  \textbf{Section 4 --- Key Takeaways}
  \begin{itemize}[leftmargin=1.2em,itemsep=2pt]
    \item Each point in a cloud carries: \textbf{3D position},
      \textbf{timestamp}, \textbf{return intensity} (reflectivity), and
      optionally a \textbf{classification label}.
    \item Standard formats: \textbf{LAS} (industry standard),
      \textbf{LAZ} (compressed LAS), \textbf{E57} (vendor-neutral).
      Proprietary formats like PCAP require sensor-specific conversion.
    \item The Velodyne VLP-16 fires \textbf{16 lasers} at 18.08\,kHz
      $\to$ up to 300,000 points/second, streamed as 1248-byte UDP
      packets every 1.3\,ms.
    \item Raw packets contain \textbf{polar measurements} (range +
      azimuth).  Converting to 3D world coordinates requires the sensor's
      factory-calibrated vertical beam angles and, for a moving platform,
      GPS/INS timestamps.
  \end{itemize}
\end{keybox}

% ==============================================================================
\newpage
\section*{Section 5 \quad Robot Position Estimation with LiDAR}
% ==============================================================================

We now apply everything from the previous sections to a concrete engineering
problem: a mobile robot equipped with a LiDAR must locate other robots in its
environment.
The robot sees a point cloud, must identify which clusters of points belong to
which robot, fit a geometric model to each cluster, and then track the robots
over time using a Kalman filter.
This section follows that pipeline step by step.

% ------------------------------------------------------------------------------
\subsection*{5.1 \quad Setup and Sensor Configuration}
% ------------------------------------------------------------------------------

\begin{figure}[H]
  \centering
  \includegraphics[width=0.82\textwidth]{slide-22.jpg}
  \caption{Mobile robots moving about in an environment, each equipped with a
    LiDAR.  The map on the right is an \emph{occupancy grid} (also called a
    cost map): each cell records whether that region of space is occupied,
    free, or unknown.  The robot uses this to plan collision-free paths.}
\end{figure}

The robots carry a \textbf{Laser Range Finder (LRF)} --- a 2D rotating LiDAR
that sweeps a single horizontal plane.
This is simpler than the full 3D VLP-16 from Section~4: rather than 16
stacked elevation angles, it fires at one fixed elevation and produces a
single ring of range measurements as it rotates.

\begin{figure}[H]
  \centering
  \includegraphics[width=0.72\textwidth]{slide-23.jpg}
  \caption{LRF sensor configuration viewed from above.  The robot body is
    the orange circle at the centre.  \textbf{Two LRF sensors} (blue dots)
    are mounted on opposite sides of the robot --- one at the front and one
    at the rear.  Each sensor sweeps its own arc (green wedges) in the scan
    direction, measuring range $d_i$ and angle $\alpha_i$ to each point it
    hits.  The small white wedge is the \textbf{blind spot}: the narrow
    gap between the two sensor fields of view where neither sensor can see,
    typically blocked by the robot body itself.}
\end{figure}

A single 2D LRF typically only covers roughly 270° of arc, leaving a
blind spot behind it.
By mounting \textbf{two sensors on opposite sides} of the robot body, the
combined coverage is nearly full 360°, with only two small blind spots
where the robot body blocks each sensor.

At each scan sample $k$, the combined sensor pair produces a
\textbf{point cloud} in polar coordinates:
\[
  S_k = \{(d_{i,k},\, \alpha_{i,k})\}, \quad i = 0, 1, \ldots, N_S
\]
where $d_{i,k}$ is the measured range and $\alpha_{i,k}$ is the angle of the
$i$-th beam at time step $k$, with measurements from both sensors merged
into one set.

% ------------------------------------------------------------------------------
\subsection*{5.2 \quad Noise Model for the Measurements}
% ------------------------------------------------------------------------------

\begin{figure}[H]
  \centering
  \includegraphics[width=0.88\textwidth]{slide-24.jpg}
  \caption{\textbf{Left:} full scan of the room.
    The \textbf{green dot} (lrf01) and \textbf{blue dot} (lrf02) mark the
    positions of the two LRF sensors mounted on the scanning robot ---
    they appear close together because both sensors sit on the same robot
    body and the room is large.
    The \textbf{red points} are all the measurements from both sensors
    combined: the large continuous cloud is the room walls; the small
    crescent-shaped clusters scattered around are the \emph{other} robots
    being detected.
    \textbf{Right:} zoomed view of the detection task.
    The scanning robot is at known position and heading $[x_0, y_0, \alpha_0]$.
    Three target robots have been detected at positions $[x_i,y_i]$,
    $[x_j,y_j]$, $[x_k,y_k]$ by fitting discs to the crescent clusters.
    The arrows show the line of sight from the scanning robot to each
    detected robot.}
\end{figure}

Every measurement is affected by noise.
The \emph{true} polar coordinates of a point are $(d_{i,k}, \alpha_{i,k})$;
the \emph{measured} (noisy) values are:
\[
  \tilde{S}_k
  = \bigl(\bar{d}_{i,k} + \nu_{d,k},\;
          \bar{\alpha}_{i,k} + \nu_{\alpha,k}\bigr)
\]
where the noise terms are independent zero-mean Gaussians:
\[
  \nu_{d,k} \sim \mathcal{N}(0,\,\sigma_{\nu_d}^2)
  \qquad
  \nu_{\alpha,k} \sim \mathcal{N}(0,\,\sigma_{\nu_\alpha}^2)
\]

\begin{notebox}
  \textbf{What the noise model tells you.}
  The range measurement $\bar{d}_{i,k}$ has Gaussian noise with standard
  deviation $\sigma_{\nu_d}$ --- a typical LiDAR might have
  $\sigma_{\nu_d} \approx 2$--$5$\,cm.
  The angle measurement $\bar{\alpha}_{i,k}$ has independent Gaussian noise
  with standard deviation $\sigma_{\nu_\alpha}$ --- typically a fraction of
  a degree.
  The two noises are \textbf{uncorrelated} (independent): an error in the
  range measurement tells you nothing about the angle error.
  This independence is captured by the diagonal noise covariance matrix
  introduced in Section~5.3.
\end{notebox}

% ------------------------------------------------------------------------------
\subsection*{5.3 \quad Uncertainty Ellipsoid: Polar to Cartesian}
% ------------------------------------------------------------------------------

\begin{figure}[H]
  \centering
  \includegraphics[width=0.72\textwidth]{slide-25.jpg}
  \caption{Uncertainty ellipsoid for a single LiDAR point.  The red dot is
    the true point $\boldsymbol{\xi}_i$; the blue dot is the noisy
    measurement $\bar{\boldsymbol{\xi}}_i$.  In polar coordinates the
    uncertainty region is a rectangle; after the nonlinear polar-to-Cartesian
    transformation it becomes a tilted ellipse in $(x,y)$ space.}
\end{figure}

\textbf{Step 1 --- Why do we need to convert to Cartesian at all?}

The LRF natively produces polar measurements $(d_i, \alpha_i)$.
Everything downstream --- clustering, disc fitting, and the Kalman filter
--- works in \textbf{Cartesian} $(x, y)$ coordinates.
So we must convert each polar point to Cartesian.
The conversion is:
\[
  \boldsymbol{\xi}_i
  = \begin{bmatrix} x_i \\ y_i \end{bmatrix}
  = \begin{bmatrix} \bar{d}_i \cos\bar{\alpha}_i \\
                    \bar{d}_i \sin\bar{\alpha}_i \end{bmatrix}
\]
This is straightforward for the \emph{position}.
The problem arises when we also want to convert the \emph{uncertainty}.

\textbf{Step 2 --- Why can we not just plug the noise directly into the
conversion formula?}

The conversion involves $\sin$ and $\cos$ --- it is \textbf{nonlinear}.
With a nonlinear function you cannot simply say ``range noise
$\sigma_{\nu_d}$ maps to Cartesian noise $\sigma_{\nu_d}$'' because the
mapping from polar error to Cartesian error changes depending on the
current angle $\alpha_i$.
A range error at $\alpha_i = 0$ moves the point purely in $x$; the same
range error at $\alpha_i = 90^\circ$ moves it purely in $y$.
The direction and magnitude of the Cartesian uncertainty depend on
\emph{where the beam is pointing} --- and you cannot capture that with a
simple scalar multiplication.

\textbf{Step 3 --- The noise covariance matrix $\Sigma$ (defined in polar).}

First, we describe the uncertainty \emph{in the polar domain}.
Since range and angle errors are independent Gaussians (from Section~5.2),
the \textbf{polar noise covariance matrix} is diagonal:
\[
  \Sigma
  = \begin{bmatrix} \sigma_{\nu_d}^2 & 0 \\
                    0 & \sigma_{\nu_\alpha}^2 \end{bmatrix}
\]
This matrix is \textbf{in polar coordinates} --- it is a $2\times2$ matrix
whose axes are range $d$ and angle $\alpha$, not $x$ and $y$.
The off-diagonal zeros reflect that range and angle errors are completely
independent: knowing one error tells you nothing about the other.
In polar space, the uncertainty region at probability level $p$ is therefore
a simple \textbf{axis-aligned rectangle}: $\pm\sigma_{\nu_d}$ in range,
$\pm\sigma_{\nu_\alpha}$ in angle.

\textbf{Step 4 --- The Jacobian linearises the conversion.}

To transform $\Sigma$ from polar into Cartesian we linearise the conversion
at the current measured point $(\bar{d}_i, \bar{\alpha}_i)$ --- exactly the
same idea as the EKF Jacobian from Lecture~4.
The Jacobian $J$ is the matrix of partial derivatives of $(x, y)$ with
respect to $(d, \alpha)$:
\[
  J(\bar{d}_i, \bar{\alpha}_i)
  = \frac{\partial\,[x,\,y]}{\partial\,[d,\,\alpha]}
  = \begin{bmatrix}
      \cos\bar{\alpha}_i & -\bar{d}_i \sin\bar{\alpha}_i \\
      \sin\bar{\alpha}_i &  \bar{d}_i \cos\bar{\alpha}_i
    \end{bmatrix}
\]
Each column tells you how a small nudge in one polar variable moves the
Cartesian point:
\begin{itemize}[leftmargin=1.2em,itemsep=2pt]
  \item \textbf{First column} ($\partial/\partial d$): a small increase in
    range moves the point in the direction $[\cos\alpha_i,\sin\alpha_i]$
    --- directly along the beam.
  \item \textbf{Second column} ($\partial/\partial\alpha$): a small increase
    in angle rotates the point perpendicular to the beam by
    $\bar{d}_i\cdot[-\sin\alpha_i,\cos\alpha_i]$ --- the $\bar{d}_i$ factor
    means this grows with range.
\end{itemize}
The Cartesian covariance is then approximated as:
\[
  \text{Cov}_\text{Cartesian} \approx J\,\Sigma\,J^\top
\]

\begin{figure}[H]
  \centering
  \includegraphics[width=0.68\textwidth]{slide-26.jpg}
  \caption{Slide~26: the uncertainty ellipse in Cartesian coordinates.
    The blue dot $\bar{\boldsymbol{\xi}}_{i,k}$ is the noisy Cartesian
    measurement; the red dot $\boldsymbol{\xi}_{i,k}$ is the true position.
    The ellipse is the set of all Cartesian points that are consistent with
    the measurement at probability level $p$.
    $\sigma_{\nu_d}$ controls the \emph{radial} (along-beam) axis;
    $\sigma_{\nu_\alpha}$ controls the \emph{tangential} (across-beam) axis,
    which grows with range because of the $\bar{d}_i$ factor in $J$.}
\end{figure}

\textbf{Step 5 --- Uncertainty ellipsoid in polar coordinates.}

We now ask: \emph{given a noisy measurement $(\bar{d}_i, \bar{\alpha}_i)$,
what is the region of true values $(d_i, \alpha_i)$ that could plausibly
have produced it?}
The answer comes from the noise model in Section~5.2: the errors
$\nu_d = \bar{d}_i - d_i$ and $\nu_\alpha = \bar{\alpha}_i - \alpha_i$
are independent Gaussians.
From statistics: the sum of $n$ independent \emph{squared standardised}
Gaussians follows a \textbf{chi-squared distribution} with $n$ degrees of
freedom.
Setting that sum equal to a threshold $\chi^2_n(1-p)$ draws a
boundary that encloses $(1-p)$ of the probability mass --- the
\textbf{confidence ellipsoid}:
\[
  \begin{bmatrix} \bar{d}_i - d_i & \bar{\alpha}_i - \alpha_i \end{bmatrix}
  \Sigma^{-1}
  \begin{bmatrix} \bar{d}_i - d_i \\ \bar{\alpha}_i - \alpha_i \end{bmatrix}
  = \chi^2_n(1-p)
\]

\textit{Expanding the matrix form.}
Because $\Sigma$ is diagonal (Step~3), so is its inverse:
\[
  \Sigma^{-1}
  = \begin{bmatrix}
      1/\sigma_{\nu_d}^2 & 0 \\
      0 & 1/\sigma_{\nu_\alpha}^2
    \end{bmatrix}
\]
Multiplying out the quadratic form gives a plain scalar equation:
\[
  \frac{(\bar{d}_i - d_i)^2}{\sigma_{\nu_d}^2}
  \;+\;
  \frac{(\bar{\alpha}_i - \alpha_i)^2}{\sigma_{\nu_\alpha}^2}
  \;=\; \chi^2_n(1-p)
\]
This is exactly the standard ellipse equation
$\dfrac{X^2}{a^2}+\dfrac{Y^2}{b^2}=1$, with semi-axis lengths:
\[
  a = \sigma_{\nu_d}\sqrt{\chi^2_n(1-p)}, \qquad
  b = \sigma_{\nu_\alpha}\sqrt{\chi^2_n(1-p)}
\]
Because there is \emph{no cross-term} $(\bar{d}_i-d_i)(\bar{\alpha}_i-\alpha_i)$
the two axes point exactly along $d$ and $\alpha$: the ellipse is
\textbf{axis-aligned} --- zero tilt.
This makes sense: range and angle errors are independent (Step~3),
so they do not ``push'' each other sideways.

\begin{notebox}
  \textbf{Understanding the chi-squared threshold $\chi^2_n(1-p)$}

  \begin{itemize}[leftmargin=1.2em,itemsep=3pt]
    \item $n = 2$: we have two polar variables ($d$ and $\alpha$), so
      $n=2$ degrees of freedom.
    \item $p$ is the probability we are willing to \emph{exclude}: choosing
      $p=0.05$ means we want a \textbf{95\,\% confidence region}.
    \item Numerical example: $\chi^2_2(0.95) \approx 5.99$.
      The semi-axis lengths become
      $\sigma_{\nu_d}\sqrt{5.99} \approx 2.45\,\sigma_{\nu_d}$
      and
      $\sigma_{\nu_\alpha}\sqrt{5.99} \approx 2.45\,\sigma_{\nu_\alpha}$.
    \item Think of it as a \emph{scaling knob}: a larger $\chi^2$ gives a
      bigger ellipse (more probability captured); a smaller $\chi^2$ gives
      a smaller ellipse (less probability but tighter bound).
    \item \textbf{Physical reading}: the ellipse is centred on the noisy
      measurement $(\bar{d}_i,\bar{\alpha}_i)$.
      The true point $(d_i, \alpha_i)$ lies inside this ellipse
      with 95\,\% probability.
  \end{itemize}
\end{notebox}

\textbf{Step 6 --- Uncertainty ellipsoid in Cartesian coordinates.}

Step~5 gives a clean axis-aligned ellipse \emph{in polar space}.
Now we want the equivalent region in Cartesian $(x,y)$, because that is
what clustering, disc fitting, and the Kalman filter all work with.

\textit{How do we get there?}
From Step~4, the Jacobian $J$ tells us how a small polar error
$\Delta m_i = [\bar{d}_i - d_i,\;\bar{\alpha}_i - \alpha_i]^\top$
maps to a Cartesian error:
\[
  \underbrace{(\bar{\boldsymbol{\xi}}_i - \boldsymbol{\xi}_i)}_{\text{Cartesian error}}
  \;\approx\;
  J \;\underbrace{[\bar{d}_i-d_i,\;\bar{\alpha}_i-\alpha_i]^\top}_{\text{polar error}}
\]
Substituting the polar error (expressed in terms of the Cartesian error)
back into Step~5's quadratic form and rearranging replaces $\Sigma^{-1}$
with the \textbf{Cartesian precision matrix} $J\Sigma^{-1}J^\top$:
\[
  (\bar{\boldsymbol{\xi}}_i - \boldsymbol{\xi}_i)^\top
  \; J \Sigma^{-1} J^\top \;
  (\bar{\boldsymbol{\xi}}_i - \boldsymbol{\xi}_i)
  = \chi^2_n(1-p)
\]
The $\chi^2_n(1-p)$ threshold is \textbf{unchanged} ---
changing coordinate systems does not change probability levels.

\textit{Why is this ellipse now tilted instead of axis-aligned?}
In Step~5, $\Sigma^{-1}$ was diagonal: its axes pointed along $d$ and
$\alpha$ with zero tilt.
After multiplying by $J$ and $J^\top$, the matrix $J\Sigma^{-1}J^\top$
is \textbf{no longer diagonal}: it contains $\sin\bar{\alpha}_i$ and
$\cos\bar{\alpha}_i$ (evaluated at the actual beam angle), which
\emph{mixes} the range and angle uncertainties.
The natural axes of this new matrix are:
\begin{itemize}[leftmargin=1.2em,itemsep=6pt]
  \item \textbf{Radial axis} --- along the beam, direction
    $[\cos\bar{\alpha}_i,\,\sin\bar{\alpha}_i]$.
    This axis captures the \emph{range} error $\sigma_{\nu_d}$.
    Range noise is already measured in \emph{metres}: if the sensor
    reads the distance 3\,cm too long, the point moves 3\,cm directly
    along the beam --- regardless of whether the object is 2\,m or
    20\,m away.
    So the half-length of this axis is simply $\approx\sigma_{\nu_d}\sqrt{\chi^2}$
    --- \textbf{constant, does not grow with range}.

  \item \textbf{Tangential axis} --- perpendicular to the beam, direction
    $[-\sin\bar{\alpha}_i,\,\cos\bar{\alpha}_i]$.
    This axis captures the \emph{angle} error $\sigma_{\nu_\alpha}$.
    Angle noise is measured in \emph{radians} --- but what we care about
    is the sideways position error in \emph{metres}.
    Think of the object sitting at the tip of the laser beam: if the
    beam swings by a small angle $\sigma_{\nu_\alpha}$, the tip traces
    an arc of length:
    \[
      \text{arc} = \bar{d}_i \times \sigma_{\nu_\alpha}
    \]
    (arc length = radius $\times$ angle).
    At 2\,m range a $0.01$\,rad error moves the point $0.02$\,m
    sideways; at 20\,m the same $0.01$\,rad moves it $0.20$\,m
    sideways.
    So the half-length is $\approx\bar{d}_i\,\sigma_{\nu_\alpha}\sqrt{\chi^2}$
    --- \textbf{grows linearly with range}.
\end{itemize}
Unless the beam points exactly along the $x$- or $y$-axis, neither of
these directions is horizontal or vertical --- the ellipse is therefore
\textbf{tilted} by exactly the beam angle $\bar{\alpha}_i$.

\begin{notebox}
  \textbf{Comparing the two ellipsoids side by side}

  \medskip
  \begin{tabular}{@{} p{2.6cm} p{4.8cm} p{5.0cm} @{}}
    \toprule
    & \textbf{Polar (Step 5)} & \textbf{Cartesian (Step 6)} \\
    \midrule
    Error vector
      & $[\bar{d}_i-d_i,\;\bar{\alpha}_i-\alpha_i]^\top$
      & $\bar{\boldsymbol{\xi}}_i - \boldsymbol{\xi}_i$ \\[4pt]
    Weight matrix
      & $\Sigma^{-1}$ (diagonal)
      & $J\Sigma^{-1}J^\top$ (full, off-diagonal) \\[4pt]
    Ellipse tilt
      & Axis-aligned along $d$, $\alpha$
      & Tilted by beam angle $\bar{\alpha}_i$ \\[4pt]
    Radial half-axis
      & $\sigma_{\nu_d}\sqrt{\chi^2}$
      & $\sigma_{\nu_d}\sqrt{\chi^2}$ (same) \\[4pt]
    Tangential half-axis
      & $\sigma_{\nu_\alpha}\sqrt{\chi^2}$
      & $\bar{d}_i\,\sigma_{\nu_\alpha}\sqrt{\chi^2}$ \emph{(grows!)} \\[4pt]
    \bottomrule
  \end{tabular}

  \medskip
  \textbf{Why does the tangential axis grow?}
  An angular error of 1\,radian at range $\bar{d}_i$ sweeps an arc of
  length $\bar{d}_i \times 1\,\text{rad} = \bar{d}_i$ metres.
  So the same angle noise produces a \emph{larger sideways position
  error} the further away the object is.
  At short range the ellipse is compact in all directions (you know
  both distance and lateral position well).
  At long range it becomes very elongated sideways: you know the
  distance to the object accurately, but its lateral position is
  increasingly uncertain --- the tangential axis keeps growing while
  the radial axis stays fixed.
\end{notebox}

% ------------------------------------------------------------------------------
\subsection*{5.4 \quad Recognising Objects: Clustering and Disc Fitting}
% ------------------------------------------------------------------------------

\begin{figure}[H]
  \centering
  \includegraphics[width=0.88\textwidth]{slide-27.jpg}
  \caption{DBSCAN clustering steps: (1) select a random point, (2) group
    nearby points within radius $\varepsilon$, (3) merge overlapping groups,
    repeat until all points are assigned.  The result (4) is a set of
    labelled clusters, one per detected object.}
\end{figure}

Given the raw point cloud $S_k$, the first task is to find which points
belong to which robot.
This involves three steps:

\textbf{Step 1 --- Clustering with DBSCAN.}

\textbf{DBSCAN} (Density-Based Spatial Clustering of Applications with
Noise) groups points by how densely packed they are.
It has two parameters:
\begin{itemize}[leftmargin=1.2em,itemsep=2pt]
  \item $\varepsilon$ --- the \emph{neighbourhood radius}: how close two
    points must be to be considered neighbours.
  \item \textsc{MinPts} --- the \emph{minimum neighbour count}: how many
    neighbours a point needs to be able to ``grow'' a cluster.
\end{itemize}

\textit{How it works, step by step:}
\begin{enumerate}[leftmargin=1.6em,itemsep=3pt]
  \item Pick any unvisited point $P$.
  \item Count how many other points lie within radius $\varepsilon$ of
    $P$ --- call this the \emph{neighbourhood} of $P$.
  \item If the neighbourhood contains \emph{fewer} than \textsc{MinPts}
    points: mark $P$ as \textbf{noise} for now (it may still be rescued
    later by another cluster passing through it).
  \item If the neighbourhood contains \emph{at least} \textsc{MinPts}
    points: $P$ is a \textbf{core point} --- start a new cluster and
    add $P$ plus all its neighbours to it.
  \item For every new point $Q$ just added to the cluster: look at
    $Q$'s own neighbourhood.
    If $Q$ is also a core point ($\geq$\,\textsc{MinPts} neighbours),
    pull \emph{its} neighbours into the cluster too.
    Keep doing this until no new points can be added --- the cluster
    has been fully expanded.
  \item Move to the next unvisited point and repeat from step~1.
\end{enumerate}

After all points are visited, every point is either:
\begin{itemize}[leftmargin=1.2em,itemsep=2pt]
  \item \textbf{Core point}: $\geq$\,\textsc{MinPts} neighbours, can expand
    the cluster.
  \item \textbf{Border point}: $<$\,\textsc{MinPts} neighbours, but
    within $\varepsilon$ of a core point --- part of the cluster, just
    can't grow it further.
  \item \textbf{Noise}: not within $\varepsilon$ of any core point ---
    excluded entirely.
\end{itemize}

\begin{notebox}
  \textbf{Why DBSCAN works well for LiDAR point clouds}

  Each robot in the scan reflects several LiDAR points on its surface
  (an arc of points), all close together.
  Between two robots there is empty space --- no LiDAR returns.
  DBSCAN exploits exactly this:
  \begin{itemize}[leftmargin=1.2em,itemsep=2pt]
    \item Points on the \emph{same robot} are close (within $\varepsilon$)
      $\rightarrow$ they join the same cluster.
    \item Points on \emph{different robots} are separated by a gap larger
      than $\varepsilon$ $\rightarrow$ they form separate clusters.
    \item A single stray reflection has no neighbours
      $\rightarrow$ labelled noise and ignored.
  \end{itemize}
  The result is one cluster per robot --- no need to know in advance how
  many robots there are (unlike k-means, which requires you to specify
  the number of groups).
\end{notebox}
\clearpage
\subsection{Disc Fitting}
\begin{figure}[H]
  \centering
  \includegraphics[width=0.72\textwidth]{slide-28.jpg}
  \caption{Disc model for a robot cluster.  Given $N_K$ noisy 2D points
    $\xi_{i,k}$ belonging to cluster $K$, find the centre
    $\xi_{c,k} = (x_{c,k}, y_{c,k})$ and radius $\rho$ of the best-fitting
    disc.}
\end{figure}

Each robot is modelled as a \textbf{disc} (circle in 2D).
Given $N_K$ noisy Cartesian points $\boldsymbol{\xi}_{i,k}$
assigned to cluster $K$, the best-fitting disc centre
$\boldsymbol{\xi}_{c,k}$ and radius $\rho$ are found by minimising:
\[
  \min_{\boldsymbol{\xi}_{c,k},\,\rho}
  \sum_{i=0}^{N_K}
  \Bigl[
    (\boldsymbol{\xi}_{c,k} - \boldsymbol{\xi}_{i,k})^\top
    (\boldsymbol{\xi}_{c,k} - \boldsymbol{\xi}_{i,k})
    - \rho^2
  \Bigr]^2
\]
This is a least-squares problem: minimise the sum of squared differences
between the squared distances from each point to the candidate centre and
the squared radius.

\begin{notebox}
  \textbf{What are we trying to minimise, and what do we get?}

  We are trying to find the \textbf{disc centre $(x_c, y_c)$ and radius
  $\rho$} that best explain all the noisy points in the cluster.
  The cost function penalises every point that does \emph{not} sit
  exactly on the circle: if a point is too far from the centre (further
  than $\rho$) or too close, it adds to the error.
  Minimising the total error over all $N_K$ points gives the
  \textbf{best-fit circle} through the cluster.

  What we get out: the disc centre $(x_{c,k}, y_{c,k})$, which is the
  \textbf{estimated robot position} at time step $k$.
  This is then passed to the Kalman filter for tracking.
\end{notebox}

\textbf{Reformulating as a linear problem.}

\begin{figure}[H]
  \centering
  \includegraphics[width=0.82\textwidth]{slide-29.jpg}
  \caption{Reformulation of the disc-fitting problem as a linear least-squares
    system $\min_p \|X^\top p - q\|_2$.  The matrix $X$ and vector $q$ are
    built from the measured point coordinates; $p$ contains the unknowns
    (centre and radius).}
\end{figure}

\textit{Why is the original problem nonlinear?}
A perfect circle satisfies $(x_i - x_c)^2 + (y_i - y_c)^2 = \rho^2$ for
every point $i$.
The unknowns $x_c$, $y_c$ sit \emph{inside} squared terms, so minimising
the error directly requires iterative nonlinear solvers.
The trick is to expand the bracket algebraically and rearrange so that
the unknowns appear only \emph{multiplied by known numbers} (i.e.\
linearly).

\textit{The algebraic trick --- expand and rearrange.}
Start from the ideal circle equation for point $i$:
\[
  (x_{i,k} - x_{c,k})^2 + (y_{i,k} - y_{c,k})^2 = \rho^2
\]
Expand:
\[
  x_{i,k}^2 - 2x_{i,k}x_{c,k} + x_{c,k}^2
  + y_{i,k}^2 - 2y_{i,k}y_{c,k} + y_{c,k}^2 = \rho^2
\]
Move all the \emph{known} data terms ($x_{i,k}^2 + y_{i,k}^2$) to the
right and all \emph{unknown} terms to the left:
\[
  2x_{c,k}\cdot x_{i,k}
  + 2y_{c,k}\cdot y_{i,k}
  + \underbrace{(\rho^2 - x_{c,k}^2 - y_{c,k}^2)}_{\text{one combined unknown}}
  \cdot 1
  = x_{i,k}^2 + y_{i,k}^2
\]
Now define three new unknowns that group the messy parts together:
\[
  p_1 = 2x_{c,k}, \qquad
  p_2 = 2y_{c,k}, \qquad
  p_3 = \rho^2 - x_{c,k}^2 - y_{c,k}^2
\]
The equation for point $i$ becomes:
\[
  p_1 \cdot x_{i,k} \;+\; p_2 \cdot y_{i,k} \;+\; p_3 \cdot 1
  \;=\; x_{i,k}^2 + y_{i,k}^2
\]
This is \textbf{linear} in $(p_1, p_2, p_3)$ --- all the unknowns appear
only multiplied by known numbers.

\textit{Writing all $N_K$ equations as one matrix system.}
We have one equation per point.
Stack them all into a matrix form $X^\top p = q$:
\[
  \underbrace{
  \begin{bmatrix}
    x_{1,k} & y_{1,k} & 1 \\
    x_{2,k} & y_{2,k} & 1 \\
    \vdots  & \vdots  & \vdots \\
    x_{N_K,k} & y_{N_K,k} & 1
  \end{bmatrix}
  }_{X^\top \;(N_K \times 3)}
  \underbrace{
  \begin{bmatrix} p_1 \\ p_2 \\ p_3 \end{bmatrix}
  }_{p\;(3\times 1)}
  =
  \underbrace{
  \begin{bmatrix}
    x_{1,k}^2 + y_{1,k}^2 \\
    x_{2,k}^2 + y_{2,k}^2 \\
    \vdots \\
    x_{N_K,k}^2 + y_{N_K,k}^2
  \end{bmatrix}
  }_{q\;(N_K \times 1)}
\]

\begin{notebox}
  \textbf{What each matrix and vector actually contains}

  \medskip
  \textbf{$X^\top$ ($N_K \times 3$) --- one row per point:}
  \begin{itemize}[leftmargin=1.2em,itemsep=2pt]
    \item Column 1: the $x$-coordinate of each Cartesian point
      (will be multiplied by $p_1 = 2x_c$).
    \item Column 2: the $y$-coordinate of each Cartesian point
      (will be multiplied by $p_2 = 2y_c$).
    \item Column 3: all ones (will be multiplied by $p_3$, the
      combined constant term).
  \end{itemize}
  Equivalently, $X$ is $3 \times N_K$ with each \emph{column} being
  one point $[x_{i,k},\, y_{i,k},\, 1]^\top$.

  \medskip
  \textbf{$q$ ($N_K \times 1$) --- one entry per point:}
  Each entry is $x_{i,k}^2 + y_{i,k}^2$ --- the \emph{squared distance
  from the origin} to point $i$.
  This is the right-hand side that came from moving the data terms
  across in the rearrangement above.
  It is built entirely from known measured values.

  \medskip
  \textbf{$p$ ($3 \times 1$) --- the three unknowns:}
  \begin{itemize}[leftmargin=1.2em,itemsep=2pt]
    \item $p_1 = 2x_{c,k}$ \quad$\Rightarrow$ extract centre
      $x_{c,k} = p_1/2$
    \item $p_2 = 2y_{c,k}$ \quad$\Rightarrow$ extract centre
      $y_{c,k} = p_2/2$
    \item $p_3 = \rho^2 - x_{c,k}^2 - y_{c,k}^2$
      \quad$\Rightarrow$ extract radius
      $\rho = \sqrt{p_3 + x_{c,k}^2 + y_{c,k}^2}$
  \end{itemize}
\end{notebox}

\textit{Why is the system overdetermined, and what does the normal
equation do?}

We have $N_K$ equations (one per point) but only 3 unknowns.
As long as $N_K > 3$ --- which is always true for a useful cluster ---
there is no vector $p$ that satisfies all equations exactly, because
the measured points are noisy.
Instead we \textbf{minimise the total squared residual}
$\|X^\top p - q\|^2$: find the $p$ whose predicted values
$X^\top p$ are as close as possible to the measured right-hand side $q$.

Differentiating $\|X^\top p - q\|^2$ with respect to $p$ and setting
it to zero gives the \textbf{normal equation}:
\[
  XX^\top\, p = Xq
  \quad\Longrightarrow\quad
  p = (XX^\top)^{-1}\, Xq
\]
Note that $XX^\top$ is only $3\times 3$ (from a $3\times N_K$ times
$N_K\times 3$ product), regardless of how many points there are ---
so the matrix inversion is tiny.
This gives the \textbf{exact global minimum in one step}: no
iteration, no initial guess needed.

\begin{keybox}
  \textbf{What do you actually get out of this minimisation?}

  You get the disc centre $(x_{c,k},\, y_{c,k})$ --- the
  \textbf{best estimate of the robot's position} given all $N_K$ noisy
  measurements in that cluster.

  Think of it this way:
  \begin{itemize}[leftmargin=1.2em,itemsep=3pt]
    \item Each individual LiDAR point on the robot's surface is
      \emph{noisy} --- it has its own uncertainty ellipsoid from
      Section~5.3.
      You could use one point as your position guess, but it would be
      unreliable.
    \item The disc fitting uses \emph{all} $N_K$ points at once.
      The noise on different points is independent --- some measure a bit
      too far, some a bit too close, some shifted left, some shifted right.
      By fitting a circle through all of them together, those random
      errors largely \emph{cancel out}, and the centre estimate becomes
      much more reliable than any single point alone.
    \item The more points in the cluster, the better: noise averages
      out faster with more measurements.
  \end{itemize}
  So the result is not truly ``noise-free'' --- it is the
  \textbf{maximum likelihood estimate} of the disc centre: the most
  probable true robot position consistent with all the noisy data.
  This estimated centre $(x_{c,k}, y_{c,k})$ is then handed to the
  \textbf{Kalman filter} as the position measurement for that time step.
\end{keybox}

% ------------------------------------------------------------------------------
\subsection*{5.5 \quad Simulation and Tracking Moving Targets}
% ------------------------------------------------------------------------------

\begin{figure}[H]
  \centering
  \includegraphics[width=0.80\textwidth]{slide-30.jpg}
  \caption{Simulation result with the LiDAR placed at the origin $(0,0)$.
    Three robots (blue, green, red) are detected from partial arc-shaped
    point clusters on their surfaces.  The fitted disc centres ($+$ markers)
    closely match the true robot positions.  Note that each robot only
    reflects a partial arc --- the LiDAR only sees the side facing the
    sensor.}
\end{figure}

Once the disc centre $\boldsymbol{\xi}_{c,k} = (x_{c,k}, y_{c,k})$ of each
robot is found at each time step $k$, the position estimate can be fed into
a \textbf{Kalman filter} to track the robots as they move.

\begin{figure}[H]
  \centering
  \includegraphics[width=0.88\textwidth]{slide-31.jpg}
  \caption{State-space model for a moving robot.  The state vector contains
    position and velocity; $\Phi$ propagates the state forward by one time
    step $T_s$; $H$ extracts the position measurements from the state.}
\end{figure}

\textbf{State-space model.}
If no information about the robots' dynamics is available, a simple
constant-velocity model is used.
The state vector for one robot is:
\[
  \chi_k = \begin{bmatrix} x_k \\ y_k \\ v_{x,k} \\ v_{y,k} \end{bmatrix}
\]
containing position $(x_k, y_k)$ and velocity $(v_{x,k}, v_{y,k})$.

The state transition matrix (assuming constant velocity over one step $T_s$):
\[
  \Phi = \begin{bmatrix}
    1 & 0 & T_s & 0  \\
    0 & 1 & 0   & T_s \\
    0 & 0 & 1   & 0  \\
    0 & 0 & 0   & 1
  \end{bmatrix}
\]
Each position updates by adding velocity $\times$ time step; velocity is
assumed constant (changed only by process noise $w_k$).

The measurement matrix extracts only the position from the state, since
the LiDAR measures position (disc centre), not velocity:
\[
  H = \begin{bmatrix} 1 & 0 & 0 & 0 \\ 0 & 1 & 0 & 0 \end{bmatrix}
\]

The full discrete-time model is:
\[
  \chi_{k+1} = \Phi\,\chi_k + w_k,
  \qquad
  \xi_k = H\,\chi_k + \nu_k
\]
where $w_k$ and $\nu_k$ are Gaussian noise signals.
A standard Kalman filter can now be applied to estimate position and
velocity from the noisy disc-centre measurements.

\begin{notebox}
  \textbf{Why use a Kalman filter here?}

  The disc-fitting step gives a noisy position estimate
  $\boldsymbol{\xi}_{c,k}$ at each time step.
  On its own, each estimate has uncertainty from the LiDAR noise (via the
  uncertainty ellipsoid from Section~5.3).
  The Kalman filter does two things:
  \begin{itemize}[leftmargin=1.2em,itemsep=2pt]
    \item \textbf{Smooths} the position estimate by combining the current
      measurement with the prediction from the dynamic model --- a single
      noisy measurement is less trusted than the combination of many
      measurements over time.
    \item \textbf{Estimates velocity} even though velocity is never directly
      measured --- the filter infers it from how the position changes between
      steps.
  \end{itemize}
  For multiple robots, independent copies of the state vector and filter
  are run in parallel --- one per robot.
\end{notebox}

% ── Key takeaways ─────────────────────────────────────────────────────────────
\begin{keybox}
  \textbf{Section 5 --- Key Takeaways}
  \begin{itemize}[leftmargin=1.2em,itemsep=2pt]
    \item The LRF produces a polar point cloud
      $S_k = \{(d_{i,k}, \alpha_{i,k})\}$ at each time step.
      Noise is modelled as independent Gaussians on range and angle.
    \item Polar $\to$ Cartesian uses a nonlinear transform; the
      \textbf{Jacobian} $J$ linearises the noise propagation.
      The result is a tilted ellipse whose tangential axis grows with
      range: $\sigma_\text{tangential} \approx \bar{d}_i\,\sigma_{\nu_\alpha}$.
    \item \textbf{DBSCAN} clusters the point cloud into groups, one per
      object.  Each cluster is fitted with a \textbf{disc} (circle) model
      via linear least-squares: $p = (XX^\top)^{-1}Xq$, giving centre
      and radius in one matrix step.
    \item The disc centre positions are fed into a \textbf{Kalman filter}
      with a constant-velocity state model
      $\chi_k = [x, y, v_x, v_y]^\top$ to track moving robots and
      estimate velocity from position measurements alone.
    \item For multiple robots, independent Kalman filter instances run
      in parallel --- one state vector per robot.
  \end{itemize}
\end{keybox}

\end{document}
